\chapter{Ley General de Masas: escala absoluta de acción \texorpdfstring{\(10^{-34}\)}{10⁻³⁴}, dominio \texorpdfstring{\(81/56/13/324\)}{81/56/13/324}, sector nulo y construcción operatorial del espectro material}
\label{chap:ley-general-masas-dominio-completo}

\section{Dominio operatorial de la Ley General de Masas: 81 celdas, 56 familias, 13 multisecciones, 324 rutas y sector nulo}
La Ley General de Masas actúa sobre secciones persistentes del atlas
orientado, no sobre una tabla de especies. Su dominio interno contiene
\(81\) celdas, \(56\) familias de trayectoria, \(13\)
multisecciones y \(324\) rutas orientadas, además de una
carta nula separada. Los censos de salidas físicas son posteriores y no
alteran este dominio.

\section{Origen determinantal de la escala absoluta}
La forma entera \(H_4\) tiene forma normal de Smith
\(\operatorname{diag}(1,2,2,4)\), por lo que su conúcleo posee orden \(16\).
De aquí se obtiene
\[
 \Sigma_H=\varphi_{\rm HMT}\alpha_{\rm HMT}^{16},
 \qquad \operatorname{ord}_{10}(\Sigma_H)=-34.
\]
La cadena absoluta es
\[
 \Sigma_H\to\hbar_{\rm pre/ret}\to\omega_0
 \to\varepsilon_{\rm clk}\to m_{\rm clk}
 \to\widehat M^{(0)}\to\widehat M^{(\beta)}.
\]
La década \(-34\) nunca se introduce desde el sistema SI ni se absorbe en el
corrector relativo.

\section{Construcción rectora de la ley y del dominio material}
El desarrollo integral de la ley se distribuye aquí por resultado: estado, ruta,
registro, firma, carácter, torres, rama absoluta y operador. Su orden interno
se conserva, mientras las realizaciones por especie se reservan al capítulo
correspondiente.

\begingroup
\let\chapter\section
\let\section\subsection
\let\subsection\subsubsection
\let\subsubsection\paragraph
\let\clearpage\relax
\section{Definición espectral de la masa como medida de persistencia de \(1\) ruta acoplada a un régimen físico}
La masa es la medida espectral de la persistencia de una ruta cuando una sección genealógica se acopla a un régimen físico. No es un atributo escalar adherido de antemano a un nombre de partícula. La dinámica APP--TRIT--TPK construye la ruta, su memoria, su firma y su carácter; la Mecánica Dimensional realiza esa genealogía en una recta de masa y determina, mediante el acoplamiento, el subespacio observable y su espectro.

Una partícula elemental se representa por una sección persistente; una familia no central, por una multisección estable; una partícula compuesta, por una composición de registros genealógicos con frontera de enlace; una resonancia, por un polo del operador acoplado; un sector topológico, por su anillo de fusión y su representación de trenzado antes de toda carta de masa; una sección virtual, por su horizonte terminal de prolongación. Estas realizaciones son distintas y no deben ser aplanadas en una sola lista de números.

La magnitud externa nunca interviene en la selección de celda, ruta, firma, lector, coeficiente de torre ni autovector. La confrontación metrológica se efectúa después de construir y fijar la salida interna.

\Needspace{7\baselineskip}
\section{Espacio completo de estados de masa}
Sea \(\mathcal C_{81}\) el atlas celular. Cada celda admite cuatro orientaciones two--way, por lo que el espacio finito de rutas orientadas es

\[
 \mathcal R_{324}
 =\{(c,h_0,v_0):c\in\mathcal C_{81},\ (h_0,v_0)\in\{\pm1\}^2\}.
\]
La proyección celular, la relación de familia y la estratificación gruesa producen

\[
 |\mathcal C_{81}|=81,\qquad
 |\mathcal F_{56}|=56,\qquad
 |\mathcal M_{13}|=13,\qquad
 |\mathcal R_{324}|=324.
\]
El electrón ocupa el único punto fijo central. Las doce multisecciones no centrales son fibras finitas estables bajo la inversión celular y no admiten, en general, una sección puntual equivariante. Su objeto natural es, por tanto, un operador sobre la fibra y no un número elegido entre sus celdas. El sector nulo se incorpora mediante la sección cero de la recta de masa, no como valor de un carácter positivo.

El catálogo físico de salida distingue tres leptones cargados, tres realizaciones neutrínicas de interacción, seis sabores quark con sus órbitas de color, dos realizaciones de calibre nulas, dos bosones de calibre masivos, un sector escalar, dos familias compuestas, tres codominios topológicos y un sistema de secciones virtuales. Esta clasificación posterior no limita el número de estados compuestos o resonantes: composición y excitación generan familias adicionales dentro de sus codominios respectivos.

\Needspace{7\baselineskip}
\section{Trayectoria genealógica de una realización de masa}
La ley no comienza con una etiqueta de partícula ni con un número. Su orden es

\[
\boxed{
\begin{aligned}
\mathrm{APP}\to\mathrm{TRIT}\to\mathrm{TPK}
&\to\widetilde\Gamma^{\rm surv}
\to\Gamma^{\rm obs}
\to L_{108}
\to\sigma\in\mathbb Z^5\\
&\to\chi_0
\to\chi_{\rm full}
\to\widehat{\mathcal M}
\to\text{acoplamiento}\\
&\to\text{espectro/polo}
\to\mathcal L_M
\to\text{contraste}.
\end{aligned}}
\]
Aquí:

\begin{itemize}
\item APP dispone la incompatibilidad local, las hojas suma/producto y su orientación;
\item el TRIT conserva signo, umbral y apertura bidireccional;
\item el TPK evoluciona fase, memoria, acarreo, supervivencia y retorno;
\item una partícula es una sección persistente de una extensión, y su ruta observable es
la historia que deja en el registro genealógico;

\item la firma entera es aditiva;
\item el carácter convierte composición aditiva en una razón positiva;
\item las torres refinan el carácter mediante un único operador armónico global;
\item la recta de masa aporta el tipo dimensional;
\item medir es acoplarse: el acoplamiento determina qué subespacio y qué espectro son
físicamente accesibles, sin consultar la magnitud que después se medirá.

\end{itemize}
La masa no selecciona la ruta. La ruta persistente, el registro genealógico y el acoplamiento seleccionan el objeto espectral cuya masa se lee.

\Needspace{7\baselineskip}
\subsection{Componentes conservadas del estado enriquecido del TPK en la trayectoria de masa}
El estado que alimenta la ley de masas es

\[
 X_K=(w_6,w_{30},R_{36},g_K,B_K,p_K,c_K,\varepsilon_K,\omega_K),
\]
donde se conservan palabra corta, palabra prolongada, región, fase nonádica, memoria de bloque, prefijo, acarreo, hoja y devanado. La cadena

\[
 w_6\longrightarrow w_{30}\longrightarrow R_{36}
 \longrightarrow G_9
\]
no proyecta directamente una masa: construye la identidad persistente que después se lee.

\Needspace{7\baselineskip}
\subsection{Filtración nonádica con memoria}
Desde cada estado se obtiene una fibra de refinamiento formada por \(729\) subcilindros,

\[
 \{C_{K,j}:0\le j<729\}.
\]
La poda \(G_9\) satisface

\[
 G_9(C_{K,j})\in\{0,1\},\qquad
 \sum_{j=0}^{728}G_9(C_{K,j})=1.
\]
Si \(j_K\) es el superviviente,

\[
 P_{K+1}=729P_K+j_K.
\]
Tras nueve pasos,

\[
 g(K+9)=g(K),
\]
pero cambian, en general,

\[
 (B_{K+9},P_{K+9},\operatorname{pref}_{K+9},
 \operatorname{Witt}_{K+9})
 \ne
 (B_K,P_K,\operatorname{pref}_K,\operatorname{Witt}_K).
\]
Ésta es la forma ecuacional de la persistencia: retorna la fase, no el estado. La masa puede reconocer una partícula después del retorno porque el registro genealógico conserva la diferencia.

\Needspace{7\baselineskip}
\subsection{Transducción del acarreo nonádico al estado firmado dodecafásico}
Para una coordenada temporal \(\theta\), definimos

\[
 g_0(\theta)=\theta\bmod9,\qquad
 \beta(\theta)=
 \left\lfloor\frac{\widetilde\theta}{9}\right\rfloor\bmod12.
\]
Entonces

\[
 g_0(\theta+9)=g_0(\theta),\qquad
 \beta(\theta+9)=\beta(\theta)+1\pmod{12}.
\]
Así, el retorno nonádico conserva la fase y avanza un diente de memoria. Los 108 tiempos son

\[
 108=9\cdot12,
\]
de modo que cada registro genealógico contiene un ciclo completo de las dos escalas. Las 36.864 composiciones de acarreo describen las combinaciones locales de esta memoria bidireccional.

\Needspace{7\baselineskip}
\subsection{Completación y conservación de la unidad}
La frontera base mil se descompone como

\[
 999=729+270=729+243+27,
\]
mientras la completación es

\[
 1000=729+271=729+243+27+1,
\qquad
 271=270+1.
\]
El \(+1\) es el único vértice totalmente fronterizo y realiza el acarreo que devuelve la unidad. La partición \(729/271\) distingue núcleo activo y borde mediador; no es una aproximación áurea. La misma conservación permite que una ruta cambie de bloque sin perder su identidad normalizada.

\Needspace{7\baselineskip}
\subsection{Teorema del atlas orientado completo}
Cada celda \(c\in\mathcal C_{81}\) posee cuatro orientaciones

\[
 (h_0,v_0)\in\{\pm1\}^2.
\]
El mapa

\[
 (c,h_0,v_0)\longmapsto\Gamma(c,h_0,v_0)
\]
es biyectivo sobre \(\mathcal R_{324}\). En consecuencia,

\[
 |\mathcal R_{324}|=4|\mathcal C_{81}|=324.
\]
Cada ruta posee 108 entradas de registro genealógico, de modo que el censo temporal total es

\[
 81\cdot108=8748.
\]
La cuádrupla

\[
 (\text{fila semilla},\text{columna semilla},h_0,v_0)
\]
determina de manera única la ruta, su palabra \(w_6\), su palabra dodecafásica, sus retornos, memoria, acarreo, firma histórica y lectores. Por tanto, las 324 rutas no sólo pueden enumerarse: pueden reconstruirse y unirse exactamente, campo por campo, sin utilizar una masa.

\Needspace{7\baselineskip}
\subsection{Cocientes de familia y multisección}
El cociente de las 81 celdas por historia de ruta produce 56 familias. Sus multiplicidades se distribuyen como

\[
 41\times1,\quad10\times2,\quad2\times3,\quad1\times4,\quad2\times5,
\]
y

\[
 41+20+6+4+10=81.
\]
La estratificación en trece multisecciones posee rangos

\[
 1,\quad2,\quad4,\quad6,\quad8,\quad12,\quad16
\]
con multiplicidades

\[
 1,\quad2,\quad3,\quad2,\quad3,\quad1,\quad1,
\]
y suma de rangos

\[
 1+4+12+12+24+12+16=81.
\]
El electrón ocupa el bloque de rango uno. Los demás sectores viven en bloques de rango mayor y exigen un proyector de representación, no una celda nominal.

\Needspace{7\baselineskip}
\subsection{Firma, lectores y memoria por ruta}
Para cada \(\Gamma\in\mathcal R_{324}\) se calculan:

\[
 \sigma_\Gamma\in\mathbb Z^5,\qquad
 K_D(\Gamma),\qquad K_\Omega(\Gamma),
\]
junto con primeros retornos de clase, órbita, celda y estado cinemático; separación de prefijos; deuda de acarreo; ambigüedad; tríadas estabilizadas; hoja y devanado. La familia \(K_D\) presenta 14 perfiles; \(K_\Omega\), nueve; juntos forman 40 pares y 18 coincidencias. Esta diversidad es la información fina que permite pasar de una fibra gruesa a un operador físico.

\Needspace{7\baselineskip}
\section{Del defecto APP a la firma entera}
\Needspace{7\baselineskip}
\subsection{Defecto primordial y despliegue}
En el centro APP,

\[
 B_c=9P_++5P_-,\qquad 9-5=4.
\]
El salto local primordial es 4. El defecto global \(54\) pertenece a un despliegue posterior \(2\to6\to54\); no debe retroproyectarse sobre el centro como si fuese la misma operación.

\Needspace{7\baselineskip}
\subsection{\(12\) eventos y divisor dodecafásico}
La coordenada orientada no recibe arbitrariamente un factor \(1/6\). La cadena es

\[
12\ \text{eventos}
\to s_C^{\rm tot}
\to 1+5+6
\to s_C^{(12)}=\frac{s_C^{\rm tot}}6
\to W_\pm=6n_A\pm s_C
\to\operatorname{SNF}(1,12)
\to\chi_{\rm mass}.
\]
El \(1/6\) es el cociente orientado del cierre de doce eventos. No es el residuo (-1/12), no es la pseudoinversa de una valencia y no es la incidencia hexada--octada.

\Needspace{7\baselineskip}
\subsection{Construcción de la firma entera completa desde el registro de ruta}
El registro genealógico CT--108 produce

\[
 \sigma(\Gamma)
 =\bigl(n_A,s_C,k_{\Delta_4},b_{120},b_\zeta\bigr)\in\mathbb Z^5.
\]
Las coordenadas tienen procedencias distintas:

\begin{itemize}
\item \(n_A\): duración/persistencia dominante;
\item \(s_C\): orientación neta de la bicapa, dividida por seis tras el cierre
dodecafásico;

\item \(k_{\Delta_4}\): canal de torsión mínima;
\item \(b_{120}\): coeficiente de evento del registro;
\item \(b_\zeta\): autocorrelación de paso tres.
\end{itemize}
Se define

\[
 \zeta_{270}:=135\Delta_4^2,
\]
y se mantiene separado de

\[
 s_{270}=q_-^{270}-q_+^{270}.
\]
Las 324 rutas orientadas se unen biyectivamente con sus cinco coordenadas históricas mediante la identidad canónica de ruta. Esta unión usa únicamente estructura genealógica y no consulta ninguna magnitud metrológica.

\Needspace{7\baselineskip}
\subsection{Ángulo directo, contraángulo y canales conjugados}
La coordenada angular directa procede de la constante de estructura fina ya construida:

\[
 A=1000\alpha_{\rm HMT}
 =7.297352569283800997285105472380663\ldots{}^\circ.
\]
El contraángulo no se obtiene de una masa ni del funcional de Barbero--Immirzi. Sea

\[
 D_A=\exp\!\left(-\frac{100\pi}{9}\alpha_{\rm HMT}\right),
\]
\[
 \mathcal C_\pi(\alpha_{\rm HMT})
 =169\alpha_{\rm HMT}^{6}
 +\frac{D_A\alpha_{\rm HMT}^{7}}
        {1-D_A\alpha_{\rm HMT}},
\]
\[
 y_*=\frac{\pi}{90}(H_5+\alpha_{\rm HMT})
      -\mathcal C_\pi(\alpha_{\rm HMT}),
 \qquad
 C^*=\frac{180}{\pi}y_*.
\]
Por tanto,

\[
 C^*=2.123738338968645724261951380393\ldots{}^\circ.
\]
La coordenada \(A\) es el modo angular común y \(C^*\), el modo diferencial orientado. Su paso a la carta arquimediana es

\[
 x=A_{\rm rad}=\frac{\pi A}{180},
 \qquad
 y=C^*_{\rm rad}=\frac{\pi C^*}{180}.
\]
En el álgebra real generada por la involución \(R=\operatorname{diag}(1,-1)\), con \(P_\pm=(I\pm R)/2\), el bloque biangular es

\[
 B_{A,C^*}=AI+C^*R
 =(A+C^*)P_++(A-C^*)P_-.
\]
Ángulo más contraángulo y ángulo menos contraángulo son, así, los dos autovalores de un único objeto. La exponenciación produce los canales

\[
 q_+=e^{-x-y},\qquad q_-=e^{-x+y}.
\]
La carta es invertible:

\[
 x=-\frac12\log(q_+q_-),
 \qquad
 y=\frac12\log\frac{q_-}{q_+}.
\]
La transformación \(C^*\mapsto-C^*\) conserva los proyectores e intercambia los canales. La coordenada de retorno satisface, después de la construcción regional,

\[
 \eta_{\rm ret}=\frac{C^*}{2}-\alpha_{\rm HMT}\,{\rm deg},
\]
sin constituir una segunda regla de selección de \(C^*\).

\Needspace{7\baselineskip}
\subsection{Carácter central}
Con

\[
 \Theta_0=
 \left(A_{\rm rad},\frac{C^*_{\rm rad}}6,
 \Delta_4,s_{120},\zeta_{270}\right),
\]
el carácter central es

\[
 \chi_0(\sigma)=
 \exp\!\left(
 n_AA_{\rm rad}+\frac{s_C}{6}C^*_{\rm rad}
 +k_{\Delta_4}\Delta_4+b_{120}s_{120}+b_\zeta\zeta_{270}
 \right).
\]
Por aditividad de la firma,

\[
 \chi_0(\sigma_1+\sigma_2)=\chi_0(\sigma_1)\chi_0(\sigma_2).
\]
La unidad física no entra en esta exponencial: \(\chi_0\) es un automorfismo positivo adimensional de la recta de masa.

\Needspace{7\baselineskip}
\section{Jerarquía armónica global de torres}
Con las coordenadas (x=\allowbreak{}A\_\{\textbackslash{}rm rad\}) e (y=\allowbreak{}C\textasciicircum{}*\_\{\textbackslash{}rm rad\}) ya construidas, sean

\[
 q_\pm=\exp\!\left[-\frac{\pi}{180}(A\pm C^*)\right]
 =e^{-x\mp y}.
\]
En una convención global coherente,

\[
 T=e^{-xI-yR},\qquad
 c_n=\operatorname{Tr}(T^n)=q_-^n+q_+^n,
\]
\[
 s_n=-\operatorname{Tr}(RT^n)=q_-^n-q_+^n.
\]
Equivalente y operatorialmente,

\[
 c_n=2e^{-nx}\cosh(ny),\qquad
 s_n=2e^{-nx}\sinh(ny).
\]
La involución \(C^*\mapsto-C^*\) deja fijos los proyectores y permuta los coeficientes espectrales, o, equivalentemente, los canales asociados. De aquí se obtienen

\[
 c_{m+n}=\frac{c_mc_n+s_ms_n}{2},\qquad
 s_{m+n}=\frac{s_mc_n+c_ms_n}{2},
\]
\[
 c_n^2-s_n^2=4e^{-2nx},
\]
y la recurrencia común de orden dos. Las alturas admisibles forman

\[
 \mathfrak S=\langle90,120\rangle
 =\{90,120\}\cup\{30r:r\ge6\}.
\]
Los valores \(90,120,180,210,240,270,360,420\) son testigos genealógicos, no la torre completa. Antes del cociente, la altura 360 conserva dos historias:

\[
 k[\mathfrak S_0]
 \simeq k[X_{90},Y_{120}]/(X_{90}^4-Y_{120}^3).
\]
La elevación general se escribe

\[
 \chi_{\rm full}(\sigma,\eta)
 =\chi_0(\sigma)
  \exp\!\left(\sum_{h\in\mathfrak S}\eta_hs_h\right),
\]
con

\[
 N_{120}=b_{120}+\eta_{120}.
\]
La suma numérica no borra la doble procedencia: \(b_{120}\) viene del registro y \(\eta_{120}\) del refinamiento espectral. Tampoco puede absorber \(b_\zeta\zeta_{270}\) en \(\eta_{270}s_{270}\).

\Needspace{7\baselineskip}
\section{Rama absoluta: acción, reloj y recta de masa}
El lector determinantal produce

\[
 \operatorname{SNF}(H_4)=\operatorname{diag}(1,2,2,4),\qquad
 |\operatorname{coker}H_4|=16,
\]
\[
 \Sigma_H=\varphi\alpha_{\rm HMT}^{16},\qquad
 \operatorname{ord}_{10}(\Sigma_H)=-34.
\]
Las dos secciones orientadas de una sola recta de acción son

\[
 \hbar_{\rm pre}=H_5\,10^{-34}\mathcal U_S,
 \qquad
 \hbar_{\rm ret}=\eta_{\rm ret}\,10^{-34}\mathcal U_S.
\]
Su carta aditiva fina y su carta multiplicativa son

\[
 \delta_{2w}=H_5-\eta_{\rm ret},\qquad
 R_{\rm act}=\frac{H_5}{\eta_{\rm ret}}>0.
\]
No son dos constantes de Planck. La década común fija el orden absoluto; el cociente transporta información relativa de ida/retorno.

Declarado \(t_0\in\mathcal L_T^+\), CT--108 fija

\[
 \omega_0=\frac{2\pi}{108t_0},\qquad
 \varepsilon_{\rm clk}=\hbar_{\rm ret}\omega_0,
 \qquad
 m_{\rm clk}=\varepsilon_{\rm clk}c_{\rm int}^{-2}.
\]
El carácter actúa por homotecia sobre \(\mathcal L_M\); no se convierte en \(\hbar\). En cocientes de masas, todo factor común se cancela:

\[
 \frac{m_Y}{m_X}
 =\chi_{\rm full}(\sigma_Y-\sigma_X,\eta_Y-\eta_X)
 R_{\rm act}^{\kappa_Y-\kappa_X}.
\]
Por ello \(10^{-34}\) fija la escala absoluta, pero no corrige un residual relativo. Una modificación relativa sólo puede proceder de la ruta, el reloj específico, el lector o las torres.

Queda una distinción necesaria. HMT construye una sección interna una vez declarado \(t_0\); la coordenada MeV exige una carta. Para cerrar la carta absoluta sin una convención externa debe construirse

\[
 \Gamma\longmapsto t_\Gamma
 \quad\text{o}\quad
 \Gamma\longmapsto\omega_\Gamma.
\]
La sección 10 define un lector de reloj basado en retornos y holonomía. Su normalización absoluta queda determinada cuando el operador de ruta fija una sección positiva de la recta temporal.

\Needspace{7\baselineskip}

\endgroup

La forma normal, el conúcleo de orden dieciséis, las ablaciones y la
distinción entre \(h\) y \(\hbar\) se demuestran íntegramente en
\cref{chap:escala-accion}. Aquí se emplea ese resultado ya construido para
fijar la rama absoluta de la ley, sin reiniciar su derivación.
\begingroup
\let\chapter\section
\let\section\subsection
\let\subsection\subsubsection
\let\subsubsection\paragraph
\let\clearpage\relax
\chapter{Coordenadas angulares conjugadas y descomposición espectral}
\label{chap:canales}

El paso desde el cierre aritmético de \(\alpha\) hacia la jerarquía
dimensional se realiza mediante dos coordenadas angulares. La primera procede
directamente de \(\alphaHMT\); la segunda es la fase orientada producida por
la realización analítica regional de la microfibra quíntuple. Ambas llegan a
Mecánica Dimensional ya construidas, antes de introducir la carta exponencial
que generará los momentos globales.

\section{Transferencia de las coordenadas angulares derivadas en HMT}

Las Partes I--III construyeron dependencias distintas con raíz común en
\(\alpha_{\rm HMT}\). El carácter pre-retorno procede de
\[
(\alpha_{\rm HMT},\varphi;\,9,5,4,12,54,\text{ regla graduada})
\longrightarrow H_5,
\]
mientras que la coordenada angular y el determinante proceden de
\[
\alpha_{\rm HMT}\longrightarrow
A=1000\alpha_{\rm HMT}\,{}^\circ,
\qquad
(\alpha_{\rm HMT},\pi)\longrightarrow
D_A=e^{-100\pi\alpha_{\rm HMT}/9}.
\]
La microfibra regional fija, mediante una dependencia distinta,
\(\mathcal F_\pi\to\rho_\pi=169\), y con ella
\[
\mathcal C_\pi(\alpha_{\rm HMT})
=169\alpha_{\rm HMT}^{6}
+\frac{D_A\alpha_{\rm HMT}^{7}}
       {1-D_A\alpha_{\rm HMT}}.
\]
La rama de \(H_5\) y la corrección regional confluyen sólo en el paso
siguiente:
\[
\begin{aligned}
y_*&=\frac{\pi}{90}(H_5+\alpha_{\rm HMT})
     -\mathcal C_\pi(\alpha_{\rm HMT}),\\
C^*&=\frac{180}{\pi}y_*.
\end{aligned}
\]
En adelante, \(A\) denota la coordenada angular directa y \(C^*\) la
coordenada angular regional conjugada; estas denominaciones fijan sus tipos
sin introducir una segunda regla de selección.
Por tanto, \(H_5\) no es una salida de \(A\). Aquí \(y_*\) es la coordenada
analítica en radianes y \(C^*\) su coordenada sexagesimal. Mecánica
Dimensional recibe estas salidas ya determinadas; no vuelve a definirlas.
La coordenada \(A\) no es entrada directa de la fórmula de \(C^*\): ambos
objetos se acoplan por primera vez en \((A,C^*)\mapsto q_\pm\). La igualdad
\(\alpha_{\rm HMT}=A/1000\) es únicamente el cambio de coordenada inverso.
En esta Parte conservamos las notaciones
\[
\begin{aligned}
 A&=7.297352569283800997285105472380663^\circ,\\
 \Cstar&=2.123738338968645724261951380393\ldots{}^\circ.
\end{aligned}
\]
El redondeo empleado en las tablas es
\(2.123738338968646^\circ\). El entero \(169\), el determinante \(D_A\) y
la recurrencia regional pertenecen a la demostración anterior; ninguna
comparación SI ni el funcional hexada--octada intervienen en su generación.

La
identidad angular
\[
 \etaRet=\frac{\Cstar}{2}-\alphaAngle,
 \qquad
 \alphaAngle:=\alphaHMT\,{\rm deg},
\]
define la coordenada angular adimensional posterior al retorno; no se emplea
para generar \(\Cstar\). Equivalentemente,
una vez efectuada la construcción regional,
\[
 \Cstar=2(\etaRet+\alphaAngle).
\]
Esta última igualdad es la forma inversa de la misma relación y no una
segunda regla de selección de \(C^*\).

La realización regional corrige \(H_5\) antes de producir \(\etaRet\); por
tanto, \(H_5\),
\(\etaRet\) y \(\Cstar\) tienen tipos matemáticos distintos. El símbolo
\(\hbarHMT\) se reserva para la magnitud sobre la recta de acción obtenida
sólo después de aplicar la escala \(10^{-34}\mathcal U_S\). La comparación
con una mantisa del SI pertenece al control metrológico externo y deja
invariante la construcción anterior.

\section{Fases internas y carta exponencial}

Las fases abstractas asociadas a los dos ángulos son
\[
 a=\frac A{360},\qquad c=\frac{\Cstar}{360}
 \quad\text{en }\RR/\ZZ.
\]
Antes de aplicar una exponencial real se eligen los representantes canónicos
\(\widetilde a,\widetilde c\in[0,1)\). La exponencial actúa sobre esos
representantes, no directamente sobre las clases del cociente.
La conversión a radianes se introduce únicamente al elegir la carta
arquimediana
\[
 x:=A_{\rm rad}=\frac{\pi A}{180},\qquad
 y:=\Cstarrad=\frac{\pi\Cstar}{180},
\]
\[
 q_-=e^{-x+y},\qquad q_+=e^{-x-y}.
\]
El factor \(\pi/180\) pertenece a este cambio analítico de coordenadas. El
cierre aritmético de \(\alpha\), formulado antes de esta conversión, es
independiente de dicha elección de carta angular.

\section{Álgebra involutiva central}

Consideremos la representación regular bidimensional del álgebra real
generada por una involución. En una base espectral fijamos
\[
 R=\operatorname{diag}(1,-1),\qquad I=I_2,
\]
de modo que ambos subespacios propios tienen multiplicidad uno. Sean
\[
 P_+=\frac{I+R}{2},\qquad P_-=\frac{I-R}{2}
\]
sus proyectores espectrales. El elemento central determinado por las
coordenadas angulares es
\[
 B_{A,\Cstar}=AI+\Cstar R
 =(A+\Cstar)P_++(A-\Cstar)P_-.
\]
La involución \(\Cstar\mapsto-\Cstar\) conserva los proyectores fijos
\(P_\pm\) e intercambia los dos coeficientes espectrales ---equivalentemente,
los canales asociados---. Al aplicar el exponencial se obtiene
\[
 T=e^{-xI-yR}=q_+P_++q_-P_-.
\]
Por tanto, los dos canales son los autovalores de un único elemento del
álgebra generada por \(I\) y \(R\).

\section{Invertibilidad de la carta de canales}

\begin{theorem}[Reconstrucción de las coordenadas angulares]
Para \(q_\pm>0\),
\[
 x=-\frac12\log(q_+q_-),
 \qquad
 y=\frac12\log\frac{q_-}{q_+}.
\]
En consecuencia, la aplicación
\((A,\Cstar)\mapsto(q_+,q_-)\) es inyectiva. El intercambio
\(q_+\leftrightarrow q_-\) conserva \(A\) y transforma
\(\Cstar\mapsto-\Cstar\).
\end{theorem}
\begin{proof}
La multiplicación y el cociente de las dos exponenciales dan
\(q_+q_-=e^{-2x}\) y \(q_-/q_+=e^{2y}\). Al tomar logaritmos se recuperan
\(x\) e \(y\). Como \(\pi/180\neq0\), quedan determinados ambos ángulos.
\end{proof}

\section{Identidades de adición y recurrencia}

En esta representación bidimensional, los momentos par e impar del operador
\(T\) son
\[
 c_n=\Tr(T^n)=q_-^n+q_+^n,
 \qquad
 s_n=-\Tr(RT^n)=q_-^n-q_+^n.
\]

La hipótesis de multiplicidad uno forma parte del dominio de estas fórmulas.
En una representación con multiplicidades \(m_+\) y \(m_-\), las trazas
correspondientes serían
\(m_+q_+^n+m_-q_-^n\) y
\(m_-q_-^n-m_+q_+^n\), respectivamente.

\begin{theorem}[Leyes de composición de los momentos]
Para \(m,n\geq0\),
\[
 c_{m+n}=\frac{c_mc_n+s_ms_n}{2},
 \qquad
 s_{m+n}=\frac{s_mc_n+c_ms_n}{2},
\]
\[
 c_n^2-s_n^2=4e^{-2nx},
\]
y ambas sucesiones satisfacen
\[
 X_{n+2}=c_1X_{n+1}-e^{-2x}X_n.
\]
Además,
\[
 \partial_Ac_n=-\frac{n\pi}{180}c_n,
 \quad
 \partial_As_n=-\frac{n\pi}{180}s_n,
\]
\[
 \partial_{\Cstar}c_n=\frac{n\pi}{180}s_n,
 \quad
 \partial_{\Cstar}s_n=\frac{n\pi}{180}c_n.
\]
\end{theorem}
\begin{proof}
Se expanden los productos de
\(q_-^m\pm q_+^m\) con \(q_-^n\pm q_+^n\). La identidad cuadrática resulta
de \((u+v)^2-(u-v)^2=4uv\). La recurrencia corresponde al polinomio
característico \(z^2-c_1z+q_-q_+\), y las derivadas se obtienen mediante la
regla de la cadena.
\end{proof}

\begin{corollary}[Determinación de la jerarquía completa]
Fijados \(A\) y \(\Cstar\), todas las parejas \((c_n,s_n)\) quedan
determinadas por la recurrencia. Ninguna especie física introduce un canal
adicional.
\end{corollary}

Finalmente,
\[
 s_n=2e^{-nx}\sinh(ny),
 \qquad
 c_n=2e^{-nx}\cosh(ny).
\]
La coordenada \(A\) determina el decaimiento común y \(\Cstar\) la separación
orientada. El orden decimal de la recta de acción no se introduce en estas
fórmulas. Se deduce en el \cref{chap:escala-accion} a partir del conúcleo
integral del bloque local de Hadamard. La elección posterior de una base
expresada en \(\mathrm{J\,s}\) pertenece al transductor metrológico y no
altera las identidades adimensionales.

\endgroup
\begingroup
\let\chapter\section
\let\section\subsection
\let\subsection\subsubsection
\let\subsubsection\paragraph
\let\clearpage\relax
\section{Elipse de acción, radián y doble retorno de holonomía}
\label{sec:elipse-accion-holonomia-rev8}
\index{acción!elipse}
\index{radián!sector de acción}
\index{Möbius!doble retorno}

El teorema focal anterior fija la forma angular y el lector determinantal
fija la escala de acción. Definamos
\[
 \eta=\operatorname{artanh}\frac{C^*}{A},
 \qquad
 a_S=\sqrt{2\hbar_{\rm HMT}^{\rm ret}e^\eta},
 \qquad
 b_S=\sqrt{2\hbar_{\rm HMT}^{\rm ret}e^{-\eta}}.
\]
La órbita simpléctica
\[
 Q(\theta)=a_S\cos\theta,
 \qquad
 P(\theta)=b_S\sin\theta
\]
conserva simultáneamente la razón de forma del bloque angular y el área
fijada por la recta de acción.

\begin{theorem}[Área barrida por fase]
\label{thm:area-fase-elipse-rev8}
Para todo \(\theta\),
\[
 \operatorname{Area}(\theta)
 =\frac12\int_0^\theta(QP'-PQ')\,dt
 =\hbar_{\rm HMT}^{\rm ret}\theta.
\]
Por tanto,
\[
 \operatorname{Area}(1\ {\rm rad})=\hbar_{\rm HMT}^{\rm ret},
 \qquad
 \operatorname{Area}(2\pi)=h_{\rm HMT}^{\rm ret}.
\]
\end{theorem}

\begin{proof}
Como \(QP'-PQ'=a_Sb_S\) y
\(a_Sb_S=2\hbar_{\rm HMT}^{\rm ret}\), la integral vale
\(\hbar_{\rm HMT}^{\rm ret}\theta\). La segunda identidad usa
\(h_{\rm HMT}^{\rm ret}=2\pi\hbar_{\rm HMT}^{\rm ret}\).
\end{proof}

APP aporta el área orientada; \(A,C^*\) fijan la forma y los focos; la rama
determinantal fija la escala. La vuelta primitiva encierra \(h\) y un radián
de fase barre \(\hbar\). El radián ordinario sigue siendo
\(1\,\mathrm{rad}=180^\circ/\pi\); la construcción añade su genealogía como
sector unitario de acción, no otra constante angular.

Para una sección persistente con holonomía
\(\varepsilon_\gamma\in\{+1,-1\}\), la condición
\[
 \varepsilon_\gamma e^{iJ_\gamma/\hbar}=1
\]
distingue dos cierres. En el cilindro, \(\varepsilon_\gamma=+1\) y el ciclo
primitivo transporta \(J=2\pi\hbar=h\). En la banda de Möbius,
\(\varepsilon_\gamma=-1\): una vuelta visible cambia de hoja y transporta
\(J=\pi\hbar=h/2\); la segunda restaura orientación y completa \(h\). Por
radián de la fase interna se transporta \(\hbar\) en ambos casos; el factor
\(1/2\) pertenece a la proyección sobre la vuelta visible de la base Möbius.

\begin{figure}[htbp]
\centering
\resizebox{.98\textwidth}{!}{\begin{tikzpicture}[
  font=\scriptsize,
  box/.style={draw,rounded corners=2mm,align=center,minimum height=.9cm,
    inner xsep=5pt},
  arr/.style={-{Stealth[length=2.2mm]},thick}]

  \begin{scope}[shift={(0,0)}]
    \draw[->,hmtgray] (-3.0,0)--(3.0,0) node[right]{\(u_+\)};
    \draw[->,hmtgray] (0,-2.0)--(0,2.0) node[above]{\(u_-\)};
    \draw[hmtblue,very thick] (0,0) ellipse (2.65cm and 1.45cm);
    \fill[hmtgold] (-1.85,0) circle (2pt) node[below]{\(F_-^{(\lambda)}\)};
    \fill[hmtgold] (1.85,0) circle (2pt) node[below]{\(F_+^{(\lambda)}\)};
    \node[text=hmtblue,font=\bfseries] at (0,2.55) {elipse de tasas};
    \node[align=center,text=hmtgray] at (0,-2.4)
      {\(q_+=e^{-a_\lambda^2}\), \(q_-=e^{-b_\lambda^2}\)};
  \end{scope}

  \draw[arr,hmtgold] (3.3,0)--node[above,align=center]{misma forma\\escala de acción} (6.0,0);

  \begin{scope}[shift={(9.0,0)}]
    \draw[->,hmtgray] (-2.8,0)--(2.8,0) node[right]{\(Q\)};
    \draw[->,hmtgray] (0,-2.0)--(0,2.0) node[above]{\(P\)};
    \fill[hmtlightgold] (0,0)--(2.5,0)
      arc[start angle=0,end angle=57.2958,x radius=2.5cm,y radius=1.35cm]--cycle;
    \draw[hmtviolet,very thick] (0,0) ellipse (2.5cm and 1.35cm);
    \draw[hmtgold,thick] (0,0)--(2.5,0);
    \draw[hmtgold,thick] (0,0)--({2.5*cos(57.2958)},{1.35*sin(57.2958)});
    \node[text=hmtviolet,font=\bfseries] at (0,2.55) {elipse de acción};
    \node[align=center] at (.7,.48) {\(\theta=1\)\\área \(\hbar\)};
    \node[align=center,text=hmtgray] at (0,-2.4) {vuelta \(2\pi\): área \(h\)};
  \end{scope}

  \node[box,draw=hmtblue,fill=hmtlightblue,minimum width=3.2cm] (cyl) at (3.2,-4.3)
    {cilindro \(\varepsilon=+1\)\\\(2\pi\mapsto h\)};
  \node[box,draw=hmtviolet,fill=hmtlightviolet,minimum width=4.2cm] (mob) at (9.0,-4.3)
    {Möbius \(\varepsilon=-1\)\\\(2\pi\mapsto h/2\), \(4\pi\mapsto h\)};
  \draw[arr,hmtgray] (cyl)--node[
    above,align=center,text width=1.9cm,fill=white,inner sep=1pt
  ]{cambio de\\holonomía} (mob);
  \node[align=center,text=hmtgray,font=\footnotesize] at (6.1,-5.55)
    {La fase interna transporta \(\hbar\) por radián;\\el factor dos pertenece a la vuelta visible.};
\end{tikzpicture}
}
\caption{La elipse de tasas fija la forma; su levantamiento a la recta de acción fija
el área. El cierre cilíndrico y el doble retorno de Möbius distinguen la fase
interna de la vuelta visible.}
\label{fig:elipse-accion-holonomia-rev8}
\end{figure}

\subsection{Geometría de Hilbert--Beltrami--Klein en la elipse de acción}
\label{sec:hilbert-beltrami-klein-elipse-accion}

La misma elipse que porta el área simpléctica admite una geometría proyectiva
intrínseca. Sea
\[
 \mathbb E_{\rm act}
 :=\left\{(Q,P)\in\mathbb R^2:
 \frac{Q^2}{a_S^2}+\frac{P^2}{b_S^2}<1\right\},
 \qquad
 \mathcal L_{\rm act}(Q,P)
 :=\left(\frac{Q}{a_S},\frac{P}{b_S}\right).
\]
La aplicación lineal \(\mathcal L_{\rm act}\) identifica
\(\mathbb E_{\rm act}\) con el disco unidad \(\mathbb D\). Para dos puntos
distintos \(x,y\in\mathbb E_{\rm act}\), sean \(a,b\in\partial
\mathbb E_{\rm act}\) las intersecciones de la recta \(xy\) con la frontera,
ordenadas como \(a,x,y,b\). Definimos
\begin{equation}
 d_{\rm H}^{\rm act}(x,y)
 :=\frac12\log
 \frac{\lVert y-a\rVert\,\lVert x-b\rVert}
      {\lVert x-a\rVert\,\lVert y-b\rVert}.
 \label{eq:distancia-hilbert-elipse-accion}
\end{equation}
La expresión es la razón doble de los cuatro puntos colineales y, por ello,
es independiente de la norma auxiliar usada sobre esa recta.

\begin{proposition}[Realización de Beltrami--Klein de la elipse de acción]
\label{prop:beltrami-klein-elipse-accion}
El par \((\mathbb E_{\rm act},d_{\rm H}^{\rm act})\) es isométrico, mediante
\(\mathcal L_{\rm act}\), al modelo de Beltrami--Klein del plano hiperbólico.
Sus geodésicas no orientadas son exactamente las cuerdas abiertas de la
elipse. La frontera \(\partial\mathbb E_{\rm act}\) conserva, de manera
simultánea y tipada, la órbita simpléctica
\[
 \gamma_{\rm act}(\theta)
 =(a_S\cos\theta,b_S\sin\theta),
 \qquad
 \frac12\int_0^\theta(QP'-PQ')\,dt
 =\hbar_{\rm HMT}^{\rm ret}\theta.
\]
\end{proposition}

\begin{proof}
Las transformaciones proyectivas preservan razones dobles. Como
\(\mathcal L_{\rm act}\) es una transformación lineal invertible que lleva
la elipse al disco unidad, transporta
\eqref{eq:distancia-hilbert-elipse-accion} a la distancia de Hilbert del
disco. Ésta es la métrica de Beltrami--Klein, cuyas geodésicas son sus cuerdas.
La identidad simpléctica de la frontera es el
\cref{thm:area-fase-elipse-rev8} escrita en la carta normalizada.
\end{proof}

La carta completa queda, por tanto, tipada por
\[
 \mathfrak R_{\rm act}:
 (A,C^*,\hbar_{\rm HMT}^{\rm ret})
 \longmapsto
 (\mathbb E_{\rm act},d_{\rm H}^{\rm act},
  dP\wedge dQ,\gamma_{\rm act}).
\]
La métrica proyectiva gobierna el interior; la forma simpléctica y la ley de
área gobiernan la frontera recorrida. Una realización como región física o
como dominio de matrices de covarianza se formula mediante un funtor ulterior
que conserve explícitamente estas dos estructuras; la construcción anterior
fija ya su dominio matemático y los invariantes que ese funtor debe preservar.

\endgroup
\begingroup
\let\chapter\section
\let\section\subsection
\let\subsection\subsubsection
\let\subsubsection\paragraph
\let\clearpage\relax
\subsection[Identidad autodual de Planck sobre la periodicidad de orden 108]{Identidad autodual de
Planck en la realización circular del endomorfismo temporal de orden 108}
\label{sec:identidad-autodual-planck-ct108-rev3}
\index{Planck!identidad autodual}
\index{calendario de orden 108!realización circular}
\index{Compton!lecturas reducida y completa}

La recta de acción construida en la sección precedente admite una
realización circular compatible con la periodicidad temporal de orden 108.
Esta realización coordina tres objetos que ya poseen tipos propios: la
sección de acción, la geometría de una vuelta y la transducción
gravitatoria. La coordinación resultante proporciona una identidad de
Planck interna y prepara la recta positiva sobre la que actuará después el
operador general de masa.

\subsubsection{Cartas orientadas de acción y realización circular declarada}

Sean
\[
 \mathcal L_S^+,\quad \mathcal L_T^+,\quad \mathcal L_C^+,
 \quad \mathcal L_L^+,\quad \mathcal L_M^+,\quad \mathcal L_G^+
\]
las semirrectas positivas de acción, tiempo, velocidad, longitud y masa.
Las coordenadas anterior y posterior al retorno se escriben
\begin{equation}
 \hbar_{\rm pre}=H_5\,10^{-34}\mathcal U_S,
 \qquad
 \hbar_{\rm ret}=\eta_{\rm ret}\,10^{-34}\mathcal U_S,
 \qquad
 R_{\rm act}=\frac{\hbar_{\rm pre}}{\hbar_{\rm ret}}
 =\frac{H_5}{\eta_{\rm ret}}.
 \label{eq:planck-ct108-cartas-pre-ret-rev3}
\end{equation}
El subíndice \(a\in\{\mathrm{pre},\mathrm{ret}\}\) designará cualquiera
de ambas cartas. En cada una se conserva la distinción exacta
\begin{equation}
 h_a=2\pi\hbar_a .
 \label{eq:planck-ct108-h-hbar-rev3}
\end{equation}

Fijemos una duración elemental \(t_0\in\mathcal L_T^+\), una velocidad
interna \(c_{\rm int}\in\mathcal L_C^+\) y
\(\ell_0=c_{\rm int}t_0\). La realización circular declarada de la
periodicidad temporal de orden 108 queda determinada por
\begin{equation}
 T_{108}=108t_0,
 \qquad
 \omega_{108}=\frac{2\pi}{T_{108}},
 \qquad
 R_{108}=\frac{c_{\rm int}}{\omega_{108}}
 =\frac{54}{\pi}\ell_0,
 \label{eq:planck-ct108-realizacion-circular-rev3}
\end{equation}
y por la longitud de la vuelta completa
\begin{equation}
 C_{108}=2\pi R_{108}=c_{\rm int}T_{108}=108\ell_0.
 \label{eq:planck-ct108-perimetro-rev3}
\end{equation}
Así, \(R_{108}\) es el radio de la realización y \(C_{108}\) su
perímetro. La frecuencia angular, la frecuencia cíclica y las dos
constantes de acción se relacionan mediante
\(\omega_{108}=2\pi/T_{108}\) y \(h_a=2\pi\hbar_a\).

\subsubsection{Correspondencia entre la periodicidad de orden 108 y las lecturas de Compton}

En la carta \(a\), el cuanto circular y su sección de masa son
\begin{equation}
 E_{108,a}=\hbar_a\omega_{108},
 \qquad
 m_{108,a}=\frac{E_{108,a}}{c_{\rm int}^{2}}
 =\frac{\hbar_a\omega_{108}}{c_{\rm int}^{2}}.
 \label{eq:planck-ct108-cuanto-circular-rev3}
\end{equation}

\begin{theorem}[Identidad entre la periodicidad circular y las lecturas de Compton]
\label{thm:planck-ct108-compton-rev3}
Para cada carta \(a\in\{\mathrm{pre},\mathrm{ret}\}\), las lecturas de
Compton reducida y completa satisfacen
\begin{equation}
 \boxed{
 \bar\lambda_{\rm C}(m_{108,a})
 =\frac{\hbar_a}{m_{108,a}c_{\rm int}}=R_{108},
 \qquad
 \lambda_{\rm C}(m_{108,a})
 =\frac{h_a}{m_{108,a}c_{\rm int}}=C_{108}.}
 \label{eq:planck-ct108-identidad-compton-rev3}
\end{equation}
\end{theorem}

\begin{proof}
La sustitución de
\eqref{eq:planck-ct108-cuanto-circular-rev3} produce
\[
 \frac{\hbar_a}{m_{108,a}c_{\rm int}}
 =\frac{c_{\rm int}}{\omega_{108}}=R_{108}.
\]
La segunda igualdad resulta al aplicar simultáneamente
\eqref{eq:planck-ct108-h-hbar-rev3} y
\eqref{eq:planck-ct108-perimetro-rev3}.
\end{proof}

El teorema es homogéneo en
\((\hbar_a,c_{\rm int},t_0)\): la misma periodicidad determina el radio,
el perímetro y las dos cartas de Compton. La realización circular es la
primitiva geométrica declarada; las igualdades precedentes son sus
consecuencias exactas.

\subsubsection{Selector circular y unidad interna de Planck}

La recta gravitatoria se coordina con la carta \(a\) mediante el transductor
homogéneo
\begin{equation}
 G_a(\theta)=\theta\,
 \frac{c_{\rm int}^{5}t_0^2}{\hbar_a}
 \in\mathcal L_G^+,
 \qquad \theta\in\mathbb R_{>0}.
 \label{eq:planck-ct108-transductor-rev3}
\end{equation}
Su radio gravitatorio reducido, evaluado sobre el cuanto circular, es
\begin{equation}
 r_{{\rm g},a}(\theta)
 :=\frac{G_a(\theta)m_{108,a}}{c_{\rm int}^{2}}
 =\theta\frac{\pi}{54}\ell_0.
 \label{eq:planck-ct108-radio-gravitatorio-rev3}
\end{equation}

\begin{proposition}[Selector circular de Planck]
\label{prop:selector-circular-planck-ct108-rev3}
Dentro de la realización circular
\eqref{eq:planck-ct108-realizacion-circular-rev3}, la condición
\(r_{{\rm g},a}(\theta)=R_{108}\) selecciona una única coordenada
positiva,
\begin{equation}
 \boxed{\theta_{108}^{\rm circ}=\left(\frac{54}{\pi}\right)^2.}
 \label{eq:theta-circular-planck-ct108-rev3}
\end{equation}
La coordenada es independiente de la carta orientada \(a\).
\end{proposition}

\begin{proof}
Las ecuaciones
\eqref{eq:planck-ct108-realizacion-circular-rev3} y
\eqref{eq:planck-ct108-radio-gravitatorio-rev3} reducen la condición a
\[
 \theta\frac{\pi}{54}\ell_0=\frac{54}{\pi}\ell_0.
\]
La positividad de \(\ell_0\) y la estricta monotonía de la aplicación
\(\theta\mapsto r_{{\rm g},a}(\theta)\) proporcionan la solución y su
unicidad.
\end{proof}

Sea \(G_{108,a}^{\rm circ}:=G_a(\theta_{108}^{\rm circ})\). Las unidades internas
asociadas a esta carta satisfacen
\begin{align}
 \ell_{{\rm P},a}^{\rm circ}
 &:=\sqrt{\frac{\hbar_a G_{108,a}^{\circ}}{c_{\rm int}^{3}}}
 =R_{108},
 &
 t_{{\rm P},a}^{\circ}
 &:=\sqrt{\frac{\hbar_a G_{108,a}^{\circ}}{c_{\rm int}^{5}}}
 =\omega_{108}^{-1},
 \label{eq:planck-ct108-longitud-tiempo-rev3}\\
 m_{{\rm P},a}^{\circ}
 &:=\sqrt{\frac{\hbar_a c_{\rm int}}{G_{108,a}^{\circ}}}
 =m_{108,a},
 &
 E_{{\rm P},a}^{\circ}
 &:=m_{{\rm P},a}^{\circ}c_{\rm int}^{2}=E_{108,a}.
 \label{eq:planck-ct108-masa-energia-rev3}
\end{align}
La igualdad de radios se publica, por tanto, como
\begin{equation}
 R_{108}=\bar\lambda_{\rm C}=r_{\rm g}=\ell_{\rm P}^{\rm circ},
 \qquad
 C_{108}=\lambda_{\rm C}=2\pi r_{\rm g}=2\pi\ell_{\rm P}^{\rm circ}.
 \label{eq:planck-ct108-tres-lectores-rev3}
\end{equation}

La convención de este teorema es el radio gravitatorio reducido
\(r_{\rm g}=Gm/c_{\rm int}^{2}\). El radio de Schwarzschild satisface
\(r_{\rm s}=2r_{\rm g}\); si se adopta esta segunda convención, su cruce con
\(\bar\lambda_{\rm C}\) ocurre en
\(m=m_{{\rm P},a}^{\circ}/\sqrt2\). La declaración del radio forma, por
tanto, parte de la carta y evita trasladar inadvertidamente un factor dos
entre las dos identidades.

La coordenada \(\theta_{108}^{\rm circ}\) pertenece a la especialización
circular declarada. El selector gravitatorio general, formulado sobre la
holonomía enriquecida y su respuesta espectral, conserva su residencia en
la \cref{sec:selector-gravitatorio-rev11}. De este modo, la identidad
circular fija una normalización interna de Planck y el criterio general
fija las condiciones que debe satisfacer la realización física de la
sección gravitatoria.

\subsubsection{Covariancia pre-retorno/post-retorno}

El producto de las secciones recíprocas de acción y gravedad conserva la
geometría circular:
\begin{equation}
 \frac{m_{108,\rm pre}}{m_{108,\rm ret}}=R_{\rm act},
 \qquad
 \frac{G_{108,\rm pre}^{\circ}}{G_{108,\rm ret}^{\circ}}
 =R_{\rm act}^{-1},
 \qquad
 \hbar_a G_{108,a}^{\circ}
 =\theta_{108}^{\rm circ}c_{\rm int}^{5}t_0^2.
 \label{eq:planck-ct108-covariancia-pre-ret-rev3}
\end{equation}
En particular, \(\ell_{{\rm P},\rm pre}^{\rm circ}\) y
\(\ell_{{\rm P},\rm ret}^{\rm circ}\) coinciden con \(R_{108}\). La memoria
bidireccional reside en el cociente de las masas y en su recíproco
gravitatorio; la escala geométrica reside en el producto invariante.

La reciprocidad admite una parametrización simétrica que separa la escala
común de la orientación del retorno. Defínanse
\begin{equation}
 \begin{aligned}
 \hbar_{\rm geom}
 &:=\sqrt{\hbar_{\rm pre}\hbar_{\rm ret}},
 &
 G_{108,\rm geom}^{\circ}
 &:=\sqrt{G_{108,\rm pre}^{\circ}G_{108,\rm ret}^{\circ}},
 &
 \xi_{\rm act}&:=\frac12\log R_{\rm act}.
 \end{aligned}
 \label{eq:planck-ct108-parametrizacion-simetrica-rev3}
\end{equation}
Entonces
\[
 \hbar_{\rm pre}=\hbar_{\rm geom}e^{\xi_{\rm act}},
 \qquad
 \hbar_{\rm ret}=\hbar_{\rm geom}e^{-\xi_{\rm act}},
\]
mientras las secciones gravitatorias conjugadas satisfacen
\[
 G_{108,\rm pre}^{\circ}
 =G_{108,\rm geom}^{\circ}e^{-\xi_{\rm act}},
 \qquad
 G_{108,\rm ret}^{\circ}
 =G_{108,\rm geom}^{\circ}e^{\xi_{\rm act}}.
\]
La coordenada \(\xi_{\rm act}\) conserva la orientación
pre-retorno/post-retorno; el producto \(\hbar_aG_{108,a}^{\circ}\)
conserva la escala geométrica. Esta parametrización de las cartas de acción
no introduce por sí sola velocidades distintas de propagación.

\subsubsection{Recta autodual partida de masas}

Fijada una carta \(a\), sea el álgebra partida
\[
 \mathbb D_{\rm sp}:=\mathbb R[j]/(j^2-1),
 \qquad
 P_\pm:=\frac{I\pm j}{2},
\]
y considérese su descomposición espectral
\[
 \mathbb D_{\rm sp}
 =\mathbb R P_+\oplus\mathbb R P_-,
 \qquad
 P_\pm^2=P_\pm,\quad
 P_+P_-=0,\quad
 P_++P_-=I.
\]
La conjugación partida intercambia los proyectores,
\[
 \overline{uP_++vP_-}=vP_++uP_-,
\]
y su norma es
\[
 N_{\rm sp}(uP_++vP_-):=uv.
\]

Para \(m\in\mathcal L_M^+\), defínase la sección
\begin{equation}
 Z_a(m)
 :=\frac{m_{{\rm P},a}^{\circ}}{m}P_+
   +\frac{m}{m_{{\rm P},a}^{\circ}}P_- .
 \label{eq:planck-ct108-seccion-partida-masa-rev3}
\end{equation}
La unidad \(m_{{\rm P},a}^{\circ}\) determina asimismo la involución
\begin{equation}
 \iota_{{\rm P},a}(m)
 =\frac{(m_{{\rm P},a}^{\circ})^2}{m}.
 \label{eq:planck-ct108-involucion-masas-rev3}
\end{equation}

\begin{theorem}[Recta autodual partida de masas]
\label{thm:planck-ct108-recta-partida-masa-rev3}
Para todo \(m\in\mathcal L_M^+\),
\begin{equation}
 \boxed{
 N_{\rm sp}(Z_a(m))=1,
 \qquad
 \overline{Z_a(m)}
 =Z_a\!\left(\iota_{{\rm P},a}(m)\right).}
 \label{eq:planck-ct108-norma-conjugacion-rev3}
\end{equation}
La conjugación posee un único punto fijo positivo,
\[
 m=m_{{\rm P},a}^{\circ},
 \qquad
 Z_a(m_{{\rm P},a}^{\circ})=I.
\]
\end{theorem}

\begin{proof}
La norma es el producto de las dos coordenadas recíprocas:
\[
 \frac{m_{{\rm P},a}^{\circ}}m
 \frac m{m_{{\rm P},a}^{\circ}}=1.
\]
La conjugación intercambia ambas coordenadas, lo que equivale a sustituir
\(m\) por \((m_{{\rm P},a}^{\circ})^2/m\). El punto fijo positivo satisface
\(m^2=(m_{{\rm P},a}^{\circ})^2\).
\end{proof}

Con la rapidez
\begin{equation}
 x_a(m)=\log\frac{m}{m_{{\rm P},a}^{\circ}},
 \label{eq:planck-ct108-rapidez-masa-rev3}
\end{equation}
la sección y su conjugada se escriben
\[
 Z_a(m)=e^{-x_a}P_++e^{x_a}P_-,
 \qquad
 \overline{Z_a(m)}
 =e^{x_a}P_++e^{-x_a}P_-.
\]
Sus lectores par e impar son
\begin{equation}
 Q_a(m):=\frac{e^{x_a}+e^{-x_a}}2=\cosh x_a,
 \qquad
 A_a(m):=\frac{e^{x_a}-e^{-x_a}}2=\sinh x_a.
 \label{eq:planck-ct108-lectores-par-impar-rev3}
\end{equation}
El primero conserva la distancia al punto autodual y el segundo conserva su
orientación.

Los lectores físicos recíprocos asociados a la misma coordenada son
\[
 r_{{\rm g},a}(m)=\frac{G_{108,a}^{\circ}m}{c_{\rm int}^{2}},
 \qquad
 \bar\lambda_{{\rm C},a}(m)=\frac{\hbar_a}{mc_{\rm int}},
\]
y satisfacen
\begin{equation}
 \frac{r_{{\rm g},a}(m)}{\bar\lambda_{{\rm C},a}(m)}
 =e^{2x_a(m)},
 \qquad
 \frac{\bar\lambda_{{\rm C},a}(m)}{r_{{\rm g},a}(m)}
 =e^{-2x_a(m)}.
 \label{eq:planck-ct108-lectores-exponenciales-rev3}
\end{equation}
La conjugación intercambia los dos lectores; la identidad de Planck es su
equilibrio unitario.

\begin{exactbox}[title={Identidad de Planck y clausura electrónica}]
La identidad \(m=m_{{\rm P},a}^{\circ}\) fija el origen autodual de la recta
positiva de masas. La sección electrónica fija una posición distinta de
esa misma recta mediante una genealogía propia. La
\cref{sec:ley-general-masa-accion-autonoma} construye el operador universal
que transporta escala absoluta, carácter, torres y lectura bidireccional;
la \cref{sec:cierre-electron-two-way-rev11} realiza después su reducción
central de rango uno. Así, la unidad de Planck proporciona la identidad de
la coordenada y el electrón proporciona la primera transición material
completa sobre ella.
\end{exactbox}

\begin{proposition}[Realización genealógica del centro electrónico]
\label{prop:planck-ct108-realizacion-centro-electron-rev3}
Sea
\[
 F_e^+=\mathbb R_{>0}(P_++P_-)
\]
la semirrecta positiva de la fibra central de rango uno, sea
\(\mathcal G_e\) el espacio de genealogías
electrónicas enriquecidas y sea \(\mathcal L_M^+\) la recta positiva de
masas. Un morfismo de realización
\[
 \mathfrak R_e:
 F_e^+\times\mathcal G_e\longrightarrow\mathcal L_M^+
\]
puede enviar la identidad \(I=P_++P_-\) a una masa distinta de
\(m_{{\rm P},a}^{\circ}\) sin alterar el
\cref{thm:planck-ct108-recta-partida-masa-rev3}: la identidad de Planck fija el
origen autodual de la carta, mientras la genealogía electrónica fija un
desplazamiento sobre esa carta.
\end{proposition}

\begin{proof}
Si \(\chi_e:\mathcal G_e\to\mathbb R_{>0}\) es el carácter positivo construido
por el registro electrónico, la regla
\[
 \mathfrak R_e(\lambda I,g)
 =\lambda\,m_{{\rm P},a}^{\circ}\chi_e(g)
\]
es positiva y multiplicativa en la coordenada genealógica. Para la genealogía
neutra, \(\chi_e=1\), se recupera el origen de Planck; para una genealogía con
\(\chi_e(g)\ne1\), la misma identidad interna ocupa una posición distinta de
la recta. La clausura electrónica concreta conserva su fuente en la
\cref{sec:cierre-electron-two-way-rev11}.
\end{proof}

\paragraph{Dominio de la identidad autodual de Planck: periodicidad temporal de orden \(108\), realización circular declarada y separación de la clausura electrónica}
La recta de acción pre-retorno/post-retorno y la periodicidad temporal de
orden (108) preceden a la realización circular, que se toma como primitiva
geométrica declarada. La identidad entre la periodicidad
circular y las lecturas de Compton, el
selector \(\theta_{108}^{\rm circ}\), la unidad interna de Planck y la
involución de la recta positiva son resultados exactos una vez fijada esa
realización. La selección gravitatoria general y la clausura electrónica
se construyen por separado en las secciones indicadas.

\endgroup
\begingroup
\let\chapter\section
\let\section\subsection
\let\subsection\subsubsection
\let\subsubsection\paragraph
\let\clearpage\relax
\section{La ley general antes de sus realizaciones por especie}
\label{sec:ley-general-masa-accion-autonoma}
\index{masa!ley operatorial general}
\index{acción!escala absoluta de masa}
\index{torres espectrales!ley global de masa}

La Mecánica Dimensional recibe de HMT dos ramas ya construidas. La primera
proporciona una recta de acción con dos coordenadas orientadas; la segunda
proporciona la genealogía de cada sección persistente. La masa aparece al
hacer actuar el carácter genealógico sobre una sección de la recta de masa.
Esta confluencia es anterior a cualquier caso particular y, por ello, la
escala \(10^{-34}\) pertenece a la ley general y no a una especie aislada.

\subsection{Rama absoluta: del determinante local a la recta de masa}

El lector determinantal local fija
\begin{equation}
 \Sigma_H=\varphi_{\rm HMT}\alpha_{\rm HMT}^{16},
 \qquad \varepsilon_H:=\left\lfloor\log_{10}\Sigma_H\right\rfloor=-34 .
 \label{eq:ley-masa-sigma-menos34-autonoma}
\end{equation}
La potencia dieciséis procede del conúcleo
\(
 \mathbb Z/2\oplus\mathbb Z/2\oplus\mathbb Z/4
\)
del bloque determinantal. Sobre la recta interna de acción
\(\mathcal L_S\) se obtienen las dos secciones
\begin{equation}
 \boldsymbol\hbar_{\rm pre}
 =H_5\,10^{-34}\mathcal U_S,
 \qquad
 \boldsymbol\hbar_{\rm ret}
 =\eta_{\rm ret}\,10^{-34}\mathcal U_S,
 \label{eq:ley-masa-hbar-pre-ret-autonoma}
\end{equation}
y las cartas aditiva y multiplicativa de su desdoblamiento son
\begin{equation}
 \delta_{2w}=H_5-\eta_{\rm ret},
 \qquad
 R_{\rm act}
 =\frac{\boldsymbol\hbar_{\rm pre}}
        {\boldsymbol\hbar_{\rm ret}}
 =\frac{H_5}{\eta_{\rm ret}}>0 .
 \label{eq:ley-masa-razon-accion-autonoma}
\end{equation}

Sea \(t_0\in\mathcal L_T^+\) la duración elemental del calendario TPK y
\(c_{\rm int}\in\mathcal L_C^+\) su velocidad interna. La periodicidad de
108 iteraciones del endomorfismo temporal CT108 fija la frecuencia angular
\(\omega_0=2\pi/(108t_0)\). Las secciones absoluta anterior y posterior al
retorno son, por tanto,
\begin{equation}
 \begin{aligned}
 \mathbf m_\star^{\rm pre}
 &=\boldsymbol\hbar_{\rm pre}\,\omega_0c_{\rm int}^{-2},\\
 \mathbf m_\star^{\rm ret}
 &=\boldsymbol\hbar_{\rm ret}\,\omega_0c_{\rm int}^{-2},
 \end{aligned}
 \qquad
 \frac{\mathbf m_\star^{\rm pre}}
      {\mathbf m_\star^{\rm ret}}=R_{\rm act}.
 \label{eq:secciones-absolutas-masa-autonoma}
\end{equation}
Ambas pertenecen a
\(
 \mathcal L_M=\mathcal L_S\otimes\mathcal L_T^{-1}
 \otimes\mathcal L_C^{-2}
\).
Así se separan con precisión la década, que fija la escala absoluta, y el
cociente del desdoblamiento bidireccional pre-retorno/post-retorno, que
transporta la memoria de la evolución cerrada.

En coordenadas internas positivas
\(t_0=\widehat t_0\mathcal U_T\) y
\(c_{\rm int}=\widehat c\mathcal U_C\), la reducción escalar de cualquier
sección cuya realización física ya esté fijada adopta la forma
\begin{equation}
 \mathbf m_X^\beta
 =10^{-34}\eta_{\rm ret}\,
  \frac{2\pi}{108\widehat t_0}\,\widehat c^{-2}
  a_X^{(0)}R_{\rm act}^{\kappa_X}\,\mathcal U_M .
 \label{eq:ley-escalar-general-masas-autonoma}
\end{equation}
La coordenada \(a_X^{(0)}\) procede del carácter y de la jerarquía armónica;
\(\kappa_X\) procede del lector de ruta. En particular, la década
\(-34\) pertenece a la sección absoluta y no se incorpora a
\(\kappa_X\).

\subsection{Rama genealógica: registro, carácter y jerarquía armónica}

Cada sección superviviente recorre la cadena
\begin{equation}
 \Omega_{\rm HMT}^{\rm surv}
 \longrightarrow\Gamma_{108}^{\rm enr}
 \longrightarrow\mathcal L_{108}
 \longrightarrow\sigma\in\mathbb Z^5
 \longrightarrow(\sigma,\eta)
 \longrightarrow\chi_{\rm full}(\sigma,\eta)\in\mathbb R_{>0}.
 \label{eq:cadena-genealogica-ley-masa-autonoma}
\end{equation}
El registro conserva el orden de los eventos; la firma conserva sus cinco
invariantes enteros; el carácter los evalúa sin borrar su procedencia. Su
prolongación espectral vive en el semigrupo completo
\begin{equation}
 \mathfrak S=\langle90,120\rangle
 =\{90,120\}\cup\{30r:r\ge6\},
 \qquad
 s_h=q_-^h-q_+^h,
 \quad c_h=q_-^h+q_+^h,
 \label{eq:semigrupo-global-masas-autonoma}
\end{equation}
y adopta la forma
\begin{equation}
 \chi_{\rm full}(\sigma,\eta)
 =\chi_0(\sigma)
  \exp\!\left(\sum_{h\in\mathfrak S}\eta_hs_h\right).
 \label{eq:caracter-completo-masas-autonoma}
\end{equation}
Las alturas \(90,120,180,210,240,270,360,420\) son ocho testigos
genealógicos especialmente visibles. La ley pertenece a todo
\(\mathfrak S\); la recurrencia de segundo orden de \((s_h,c_h)\) prolonga
los testigos a la jerarquía armónica completa.

\subsection{Operador universal sobre las \(13\) multisecciones del dominio material}

El atlas aporta \(81\) posiciones, \(56\) familias de ruta, \(13\)
multisecciones y \(324\) rutas orientadas. Para una multisección \(F\), sea
\(\mathcal H_F=\bigoplus_{\Gamma\in F}\mathbb C e_\Gamma\). El carácter
basal y las torres definen el operador positivo
\begin{equation}
 \widehat{\mathcal A}^{(0)}_F e_\Gamma
 =\chi_{\rm full}(\sigma_\Gamma,\eta_\Gamma)e_\Gamma.
 \label{eq:operador-caracter-basal-masas-autonoma}
\end{equation}
Las dos historias conservadas por TPK inducen lectores
\(r\in\{D,\Omega\}\): \(K_D\) actúa sobre el registro cronológico de
108 iteraciones del endomorfismo temporal CT108 y \(K_\Omega\), sobre la palabra
dodecafásica firmada. Si
\(\widehat K_F^r e_\Gamma=\kappa_r(\Gamma)e_\Gamma\), sea
\(F_s\subseteq F\) el conjunto de rutas de la sección o bloque físico \(s\), y
defínanse
\begin{equation}
 \mathcal H_{F,s}:=\bigoplus_{\Gamma\in F_s}\mathbb C e_\Gamma,
 \qquad
 \widehat{\mathcal A}^{(0)}_{F,s}
 :=\left.\widehat{\mathcal A}^{(0)}_F\right|_{\mathcal H_{F,s}},
 \qquad
 \widehat B_{F,s}^r
 :=\left.\widehat K_F^r\right|_{\mathcal H_{F,s}}.
 \label{eq:restricciones-bloque-ley-masas-autonoma}
\end{equation}
La base diagonal de rutas empleada por los certificados vigentes hace invariante
\(\mathcal H_{F,s}\) bajo ambos operadores y convierte
\(\widehat B_{F,s}^r\) en una restricción autoadjunta. Con estos tipos, la ley
general es
\begin{equation}
 \boxed{
 \widehat{\mathbf M}^{\beta,r}_{F,s}
 =\mathbf m_\star^{\rm ret}\otimes
  R_{\rm act}^{\widehat B_{F,s}^r/2}
  \widehat{\mathcal A}^{(0)}_{F,s}
  R_{\rm act}^{\widehat B_{F,s}^r/2},
 \qquad r\in\{D,\Omega\}.}
 \label{eq:ley-operatorial-general-masas-autonoma}
\end{equation}
La escritura simétrica preserva positividad sin presuponer conmutación. En la
base simultáneamente diagonal empleada por los
certificados vigentes se reduce a
\begin{equation}
 \widehat{\mathbf M}^{\beta,r}_F
 =\mathbf m_\star^{\rm ret}\otimes
  \widehat{\mathcal A}^{(0)}_F
  \exp\!\left((\log R_{\rm act})\widehat K_F^r\right).
 \label{eq:ley-operatorial-diagonal-masas-autonoma}
\end{equation}
La ecuación conserva, en un solo objeto, la escala absoluta, el carácter
global, toda la jerarquía de torres y el corrector bidireccional dependiente
de ruta. La carta física de cada sector actúa después sobre el espectro de
\eqref{eq:ley-operatorial-general-masas-autonoma}. El sector nulo posee su
propio morfismo a cero; los registros compuestos se concatenan antes de
evaluar la firma; los sectores topológicos publican primero sus invariantes
de fusión y trenzado.

\begin{theorem}[Covariancia y cancelación de la escala común]
\label{thm:covariancia-ley-general-masas-autonoma}
La ley~\eqref{eq:ley-operatorial-general-masas-autonoma} es positiva y
covariante bajo cambio unitario de base en cada fibra. Para dos autosecciones
\(X,Y\) realizadas en la misma carta de masa,
\begin{equation}
 \frac{m_Y}{m_X}
 =\frac{\chi_{\rm full}(\sigma_Y,\eta_Y)}
        {\chi_{\rm full}(\sigma_X,\eta_X)}
  R_{\rm act}^{\kappa_r(\Gamma_Y)-\kappa_r(\Gamma_X)}.
 \label{eq:cociente-ley-general-masas-autonoma}
\end{equation}
Los factores comunes \(10^{-34}\mathcal U_S\), \(2\pi/(108t_0)\) y
\(c_{\rm int}^{-2}\) fijan la sección absoluta y se cancelan exactamente en
el cociente.
\end{theorem}

\begin{proof}
El operador \(\widehat{\mathcal A}^{(0)}_F\) es positivo porque sus
autovalores pertenecen a \(\mathbb R_{>0}\); la exponencial de un operador
autoadjunto es positiva. La conjugación por una matriz unitaria preserva el
espectro y, por tanto, la ley. Al dividir dos autovalores realizados en la
misma carta, la sección común \(\mathbf m_\star^{\rm ret}\) se simplifica y
queda~\eqref{eq:cociente-ley-general-masas-autonoma}.
\end{proof}

La fibra electrónica central tiene rango uno y constituye el primer
corolario escalar de esta ley. Las restantes multisecciones conservan el par
de operadores, sus espectros y el morfismo físico que selecciona sabor,
color, orientación o composición. Esta jerarquía permite publicar juntos el
resultado algebraico, su realización dimensional y el contraste externo sin
confundir sus dominios.

\endgroup
\begingroup
\let\chapter\section
\let\section\subsection
\let\subsection\subsubsection
\let\subsubsection\paragraph
\let\clearpage\relax
% Vista científica exacta materializada desde una rama condicional.
% Fuente: source/demostraciones_primarias/09_06h_atlas_familias_rev6.tex; rama: HMTIncludeAtlasStructural.
% No modifica el contenido matemático de la rama de origen.
\chapter[Clasificación nonádica de familias]{Clasificación nonádica de familias de trayectoria}
\label{chap:atlas-genealogia-masas}
\index{atlas nonádico}
\index{familia de ruta}
\index{multisección}
\index{sector nulo}

El atlas no comienza con una lista de partículas ni con una tabla de masas.
Comienza con un soporte finito orientado y con los lectores que conservan la
historia de sus rutas.  Su función es responder a una pregunta anterior a
cualquier comparación metrológica: ¿qué clases de persistencia puede sostener
la geometría nonádica y qué información debe conservarse para distinguirlas?

La respuesta tiene tres cardinales, pero no tres inventarios rivales:
\[
 81\ \text{celdas},\qquad
 56\ \text{familias de ruta},\qquad
 13\ \text{multisecciones familia--nivel}.
\]
El primer número cuenta posiciones orientadas; el segundo identifica las
posiciones que publican la misma firma finita de ruta; el tercero reúne las
posiciones que comparten familia y profundidad respecto del centro.  Ninguno
de ellos cuenta masas conocidas, nombres de especies o filas de una
comparación externa.

\section{La geometría completa de las \(81\) celdas}

Sea
\[
 I_9=\{1,\ldots,9\},
 \qquad
 \mathcal C_{81}=I_9\times I_9.
\]
Las coordenadas \((r,c)\) conservan las dos direcciones de APP\@. La paridad
de cada coordenada define cuatro sectores internos:
\[
 \operatorname{fam}(r,c)=
 \begin{cases}
  \ell^- ,& r\equiv1,\ c\equiv1\pmod2,\\
  \nu ,& r\equiv1,\ c\equiv0\pmod2,\\
  d ,& r\equiv0,\ c\equiv1\pmod2,\\
  u ,& r\equiv0,\ c\equiv0\pmod2.
 \end{cases}
 \tag{A1}\label{eq:familias-paridad-atlas}
\]
La profundidad en el atlas no se añade como etiqueta independiente.  Es la
distancia de Chebyshev al centro:
\[
 G(r,c)=\max\{|r-5|,|c-5|\}.
 \tag{A2}\label{eq:generacion-chebyshev-atlas}
\]
Por tanto, familia y nivel proceden de la propia posición.  En particular,
el centro no se obtiene seleccionando una especie desde fuera: es la única
celda de profundidad cero.

\begin{theorem}[Censo de familias y generaciones del atlas]
\label{thm:censo-atlas-81}
Las aplicaciones \eqref{eq:familias-paridad-atlas} y
\eqref{eq:generacion-chebyshev-atlas} descomponen el atlas de manera
disjunta como
\[
 81=45_{\rm leptones}+36_{\rm quarks}
   =25_{\ell^-}+20_{\nu}+20_{d}+16_{u}.
 \tag{A3}\label{eq:censo-81-cuatro-familias}
\]
Al refinar por \(G\), se obtienen exactamente las trece multisecciones
\[
\begin{array}{c|c|c}
\text{familia}&\text{niveles}&\text{cardinales}\\ \hline
\ell^-&G0,G2,G4&1,8,16\\
\nu&G1,G2,G3,G4&2,4,6,8\\
d&G1,G2,G3,G4&2,4,6,8\\
u&G1,G3&4,12.
\end{array}
\tag{A4}\label{eq:trece-multisecciones-atlas}
\]
La celda \((5,5)\) es la única componente de
\(\mathfrak S_{\ell^-,0}\).
\end{theorem}

\begin{proof}
En \(I_9\) hay cinco coordenadas impares y cuatro pares.  Las cuatro clases
de paridad tienen, por ello, cardinales
\[
 5\cdot5=25,\qquad
 5\cdot4=20,\qquad
 4\cdot5=20,\qquad
 4\cdot4=16,
\]
lo que prueba \eqref{eq:censo-81-cuatro-familias}; las dos primeras
clases reúnen \(45\) celdas y las dos últimas, \(36\).

Para el refinamiento, hacemos crecer cuadrados centrados en \((5,5)\).
En la clase impar--impar hay una celda en radio cero, \(3^2-1=8\)
en el cascarón de radio dos y \(5^2-3^2=16\) en el de radio cuatro.
En la clase impar--par, los rectángulos acumulados hasta radios
\(1,2,3,4\) contienen, respectivamente,
\[
 1\cdot2=2,\qquad 3\cdot2=6,\qquad
 3\cdot4=12,\qquad5\cdot4=20
\]
celdas; sus diferencias sucesivas son \(2,4,6,8\).  La clase
par--impar da el mismo cálculo con las coordenadas intercambiadas.  En la
clase par--par aparecen \(2^2=4\) celdas en radio uno y
\(4^2-2^2=12\) en radio tres.  Estos cascarones son disjuntos y agotan
\(I_9^2\), de modo que no queda otra multisección.
\end{proof}

La palabra \emph{nivel} designa en este capítulo la profundidad combinatoria
\(G\).  La lectura física posterior conserva esa geometría, pero no la
define.  Así se evita convertir los nombres de las partículas en premisas
del censo que debe explicarlas.

\section{Inversión central, \(41\) órbitas y unicidad del centro}

La inversión del atlas es
\[
 \rho_{\rm c}(r,c)=(10-r,10-c).
 \tag{A5}\label{eq:inversion-central-atlas-rev6}
\]
Como \(10-r\) tiene la misma paridad que \(r\), la inversión conserva las
cuatro familias.  Además,
\[
 |(10-r)-5|=|r-5|,
 \qquad
 |(10-c)-5|=|c-5|,
\]
y, por tanto, conserva el nivel \(G\).  Cada una de las trece
multisecciones es \(\rho_{\rm c}\)-estable.

\begin{proposition}[Órbitas y excepción central]
\label{prop:orbitas-centro-atlas-rev6}
La acción de \(\rho_{\rm c}\) sobre \(\mathcal C_{81}\) tiene exactamente
\(41\) órbitas.  Su único punto fijo es \((5,5)\).  En consecuencia, la
multisección \(\mathfrak S_{\ell^-,0}\) admite el único selector puntual
\(\rho_{\rm c}\)-equivariante; en las otras doce, el objeto natural es la
órbita o la multisección completa.
\end{proposition}

\begin{proof}
La ecuación
\[
 (r,c)=\rho_{\rm c}(r,c)
\]
equivale a \(2r=2c=10\), luego \(r=c=5\).  Las otras ochenta celdas se
agrupan en cuarenta pares, de donde \(1+40=41\) órbitas.

Sea \(y\) una multisección no central y supongamos que un selector puntual
equivariante elige \(a(y)\in\mathfrak S_y\).  Como \(\rho_{\rm c}\) actúa
trivialmente sobre el tipo familia--nivel, la equivarianza impondría
\[
 a(y)=a(\rho_{\rm c}y)=\rho_{\rm c}a(y).
\]
El punto elegido tendría que ser fijo, pero el único punto fijo está en
\(\mathfrak S_{\ell^-,0}\).  La contradicción demuestra la afirmación.
\end{proof}

Esta proposición explica por qué una especie interna no es, en general, una
casilla aislada.  El electrón central es excepcional porque su multisección
se reduce a un punto.  En las demás familias, la orientación conjugada forma
parte de la identidad que debe conservarse.

\section{Las \(56\) familias de ruta}

Cada celda porta la sombra finita de su historia de transporte:
\[
 \mathcal R(r,c)=
 \bigl(
 n_A,\ s_C,\ k_{\Delta_4},\
 \nu_{120},\ \nu_{270},\ \tau_{12}
 \bigr)_{r,c}.
 \tag{A6}\label{eq:lector-familia-ruta-atlas}
\]
Las cinco primeras coordenadas son enteras; \(\tau_{12}\) es la palabra
dodecafásica orientada.  Dos celdas pertenecen a la misma familia de ruta
si y sólo si tienen la misma sextupla \(\mathcal R\).

\begin{theorem}[Censo de familias de ruta y memoria mínima]
\label{thm:censo-56-familias-ruta-rev6}
La imagen de \(\mathcal R\) contiene exactamente \(56\) elementos.  El
histograma de tamaños de sus fibras es
\[
 41\times1,\qquad
 10\times2,\qquad
 2\times3,\qquad
 1\times4,\qquad
 2\times5.
 \tag{A7}\label{eq:histograma-fibras-ruta}
\]
En particular,
\[
 41+10+2+1+2=56,
\]
\[
 41\cdot1+10\cdot2+2\cdot3+1\cdot4+2\cdot5=81.
 \tag{A8}\label{eq:doble-censo-rutas}
\]
Fijado el orden lexicográfico dentro de cada fibra, el rango
\[
 \kappa(r,c)\in\{0,1,2,3,4\}
\]
hace inyectiva la aplicación
\[
 (r,c)\longmapsto\bigl(\mathcal R(r,c),\kappa(r,c)\bigr).
 \tag{A9}\label{eq:lift-memoria-cinco-atlas}
\]
Cinco es el tamaño mínimo de un alfabeto uniforme capaz de efectuar este
levantamiento.
\end{theorem}

\begin{proof}
La enumeración exhaustiva de \(I_9^2\) agrupa las ochenta y una sextuplas
\eqref{eq:lector-familia-ruta-atlas} por igualdad coordenada a coordenada.
Produce los cinco tamaños de fibra de
\eqref{eq:histograma-fibras-ruta}; las dos identidades
\eqref{eq:doble-censo-rutas} comprueban simultáneamente el número de clases
y la cobertura de las ochenta y una celdas.

El rango lexicográfico distingue los elementos de cada fibra, de modo que
\eqref{eq:lift-memoria-cinco-atlas} es inyectiva.  Como existen dos fibras
de cardinal cinco, cualquier alfabeto común con menos de cinco símbolos
identificaría dos celdas de una de ellas por el principio del palomar.
\end{proof}

El rango \(\kappa\) es memoria de posición dentro de una fibra; no es
carga, masa ni color.  Su función es recuperar la celda que la firma visible
de ruta había identificado con otras.  El número \(56\) cuenta clases de la
sextupla \(\mathcal R\), mientras el número \(13\) cuenta clases de la pareja
\((\operatorname{fam},G)\).  Son cocientes diferentes del mismo atlas.

\paragraph{Cocientes celular y orbital del atlas: dominios y mapas de proyección.}
La proyección cruda CT--108 y la clasificación científica leen datos
distintos de las mismas ochenta y una celdas. Si
\(q_{\rm CT}\) conserva solamente la forma de la palabra trítica
dodecafásica, mientras \(\mathcal R\) conserva la sextupla de ruta y
\(q=(\operatorname{fam},G)\) conserva familia y profundidad, el cuadro
tipado es
\[
\begin{array}{rcll}
 \mathcal C_{81}
 &\xrightarrow{\ q_{\rm CT}\ }&
 \mathcal W^{\rm raw}_{7}
 & \text{formas crudas CT--108},\\[2mm]
 \mathcal C_{81}
 &\xrightarrow{\ \mathcal R\ }&
 \Sigma^{\rm ruta}_{56}
 & \text{familias de ruta},\\[2mm]
 \mathcal C_{81}
 &\xrightarrow{\ q\ }&
 \Sigma_{13}
 & \text{multisecciones familia--nivel}.
\end{array}
\tag{A9a}\label{eq:dos-cocientes-atlas-rev8}
\]
Las siete fibras crudas tienen multiplicidades
\[
 54,2,2,2,7,7,7,
 \qquad54+3\cdot2+3\cdot7=81.
\]
Por tanto, el censo arqueológico \(81\to7\) no sustituye el atlas
\(81/56/13\). Tampoco la escritura abreviada \(81/56/13\) afirma que exista
una flecha canónica \(56\to13\): hay fibras de \(\mathcal R\) que atraviesan
más de una pareja familia--nivel. Los mapas \(\mathcal R\) y \(q\) son
cocientes paralelos y conservan invariantes diferentes. En el sector quark,
la convención activa sigue siendo
\[
 \operatorname{col}(r,c)=r-c\pmod3;
\]
la escritura histórica \(c-r\) es su convención conjugada, no otro color ni
otro cociente del atlas.

\section{Distinción entre celda, familia de ruta y multisección}

Definamos
\[
 q:\mathcal C_{81}\longrightarrow\Sigma_{13},
 \qquad
 q(r,c)=\bigl(\operatorname{fam}(r,c),G(r,c)\bigr),
 \tag{A10}\label{eq:cociente-trece-multisecciones}
\]
y, para \(y\in\Sigma_{13}\),
\[
 \mathfrak S_y^{\rm atlas}
 =q^{-1}(y)
 =\{(r,c)\in\mathcal C_{81}:q(r,c)=y\}.
 \tag{A11}\label{eq:multiseccion-atlas-rev6}
\]
Los cuatro niveles de lectura quedan entonces separados:

\begin{enumerate}[label=(\roman*)]
 \item una \emph{celda} es una posición orientada \((r,c)\), con todos sus
 datos locales;
 \item una \emph{familia de ruta} es una fibra de \(\mathcal R\): conserva
 una misma firma finita de transporte;
 \item una \emph{multisección} es una fibra de \(q\): conserva familia,
 profundidad y la totalidad de sus orientaciones conjugadas;
 \item una \emph{representación física} es la realización posterior de una
 multisección enriquecida en un espacio de estados; no es otra partición del
 tablero ni la elección arbitraria de una de sus celdas.
\end{enumerate}

Esta distinción da contenido matemático a la palabra \emph{especie}.  El tipo
interno es la multisección.  Una partícula concreta es una sección
persistente realizada en ella.  La representación física conserva la acción,
la orientación y los observables del tipo; no retrocede para redefinir el
atlas a partir de una masa ya conocida.

\section{Color residual como orientación ternaria}

En el sector quark, formado por las \(36\) celdas con \(r\) par, definimos
\[
 \operatorname{col}(r,c)=r-c\pmod3,
 \tag{A12}\label{eq:color-residual-atlas-rev6}
\]
con el diccionario
\[
 0\longmapsto R,\qquad1\longmapsto G,\qquad2\longmapsto B.
\]

\begin{proposition}[Balance de color residual]
\label{prop:color-atlas}
Las treinta y seis celdas de quark se reparten exactamente como
\[
 36=12_R+12_G+12_B.
 \tag{A13}\label{eq:balance-color-atlas-rev6}
\]
La inversión central fija \(R\) e intercambia \(G\) y \(B\).
\end{proposition}

\begin{proof}
Fijado cualquiera de los cuatro valores pares de \(r\), al recorrer
\(c=1,\ldots,9\) cada clase módulo tres aparece tres veces.  Cada residuo
tiene, por tanto, \(4\cdot3=12\) representantes.  Además,
\[
 \operatorname{col}\bigl(\rho_{\rm c}(r,c)\bigr)
 =(10-r)-(10-c)
 =-(r-c)\pmod3.
\]
La negación modular fija la clase cero e intercambia las clases uno y dos.
\end{proof}

El trítico \(R,G,B\) es, por tanto, un observable exacto del atlas y no una
decoración añadida a una tabla de quarks.  La representación de color se
construye después sobre este soporte ternario; el censo \(12+12+12\) es su
geometría previa.

\section{De la extensión superviviente a la firma}

El atlas clasifica cortes finitos.  La identidad de una partícula requiere la
historia profunda de la que ese corte procede.  Sea
\[
 \Omega_{\rm HMT}^{\rm surv}
 =\varprojlim_m\mathcal S_m
\]
el sistema inverso de estados supervivientes.  La proyección al atlas \(\pi_{81}\) y la sombra
de ruta forman la cadena
\[
 \Omega_{\rm HMT}^{\rm surv}
 \xrightarrow{\ \operatorname{sh}_{108}\ }
 \Gamma_{108}^{\rm enr}
 \xrightarrow{\ \mathcal E\ }
 \mathcal L_{108}
 \xrightarrow{\ L_{\rm tick}\ }
 \mathbb Z^5.
 \tag{A14}\label{eq:cadena-extension-firma-atlas-rev6}
\]
Aquí \(\operatorname{sh}_{108}\) publica una ruta de \(108\) pasos sin
olvidar hoja, orientación, frontera ni predecesores; \(\mathcal E\) forma el
registro dodecafásico; y \(L_{\rm tick}\) extrae
\[
 v(x)=
 \bigl(n_A,s_C,k_{\Delta_4},\nu_{120},\nu_{270}\bigr).
 \tag{A15}\label{eq:firma-cinco-extension-rev6}
\]
Ninguna de estas tres flechas consulta una masa, una unidad o el nombre de
una partícula.

La razón para conservar el registro antes de proyectar se ve en sus doce
eventos orientados \(e_0,\ldots,e_{11}\).  Para \(0\leq i\leq5\), sea
\[
 p_i=e_i+e_{i+6},
 \qquad
 a_i=e_i-e_{i+6},
 \qquad
 s_C=\sum_{i=0}^{5}p_i,
\]
y, para \(0\leq i<5\),
\[
 M_i=p_i-p_5.
\]
Entonces
\[
 p_5=\frac{s_C-\sum_{i=0}^{4}M_i}{6},
 \qquad
 p_i=p_5+M_i,
\]
\[
 e_i=\frac{p_i+a_i}{2},
 \qquad
 e_{i+6}=\frac{p_i-a_i}{2}.
 \tag{A16}\label{eq:reconstruccion-registro-doce}
\]
La identidad
\[
 1+5+6=12
\]
es así una reconstrucción reversible: un total, cinco diferencias y seis
orientaciones recuperan los doce eventos.  Proyectar a la firma de cinco
enteros antes de componer historias eliminaría precisamente la colocación y
la orientación que permiten decidir si dos rutas pueden componerse.

El cociente \(q\) se eleva a la historia profunda mediante la proyección al atlas
\(\pi_{81}:\Omega_{\rm HMT}^{\rm surv}\to\mathcal C_{81}\):
\[
 \widetilde{\mathfrak S}_y
 =\{x\in\Omega_{\rm HMT}^{\rm surv}:q(\pi_{81}x)=y\}.
 \tag{A17}\label{eq:multiseccion-enriquecida-rev6}
\]
Esta es la multisección enriquecida.  Su sombra pertenece al atlas; su ruta
pertenece a \(\Gamma_{108}^{\rm enr}\); su memoria completa pertenece al
registro; y su evaluación se aplica sólo al final.  La genealogía no es una
nota añadida al valor: es el objeto que hace posible el valor y conserva las
otras lecturas de la misma sección.

\section{Bifurcación del sector nulo antes del carácter positivo}

El carácter de masa actúa sobre la firma entera y tiene codominio positivo:
\[
 \chi_{\rm mass}:\mathbb Z^5\longrightarrow\mathbb R_{>0}.
 \tag{A18}\label{eq:caracter-positivo-atlas-rev6}
\]
Por definición, no existe una firma \(v\) con
\(\chi_{\rm mass}(v)=0\).  El sector nulo no es, por consiguiente, una
decimocuarta multisección ni el límite degenerado de las trece anteriores.

Su soporte algebraico pertenece a la rama partida
\[
 \mathcal A_{\rm sp}
 =\mathbb R[\mathsf R]/(\mathsf R^2-1),
 \qquad
 P_\pm=\frac{1\pm\mathsf R}{2}.
\]
Todo elemento se escribe de manera única como
\[
 z=r_+P_++r_-P_-,
 \qquad
 N_{\rm sp}(z)=r_+r_-,
\]
y las dos rectas
\[
 \mathbb R P_+,\qquad\mathbb R P_-
\]
son nulas porque \(N_{\rm sp}(\lambda P_\pm)=0\).  Una sección nula es una
ruta abierta transportada en una de esas direcciones.  El interior material
requiere el acoplamiento de ambas, para el cual \(r_+r_->0\).

Así quedan tipadas dos ramas:
\[
 \text{persistencia material}
 \longrightarrow\mathbb Z^5
 \xrightarrow{\chi_{\rm mass}}\mathbb R_{>0},
\]
\[
 \text{propagación nula}
 \longrightarrow\partial_{\rm null}\mathcal C_{\rm sp}^{+}.
\]
Las realizaciones fotónica y gauge leen después la segunda rama.  Su
separación del carácter positivo impide borrar fotón y gluón por no aparecer
en una lista de masas, y evita a la vez fabricar una masa nula mediante un
carácter cuyo codominio la excluye.

\section{Sistema finito de clasificación y memoria de trayectorias}

Los tres cocientes del capítulo cumplen funciones complementarias:
\[
 \mathcal C_{81}
 \xrightarrow{\ \mathcal R\ }\mathcal R_{56},
 \qquad
 \mathcal C_{81}
 \xrightarrow{\ q\ }\Sigma_{13},
 \qquad
 \mathcal C_{81}
 \xrightarrow{\ /\langle\rho_{\rm c}\rangle\ }
 \mathcal O_{41}.
\]
El primero conserva la firma de ruta; el segundo conserva familia y nivel;
el tercero conserva conjugación central.  Junto con el color residual y el
levantamiento \(\kappa\), forman un sistema finito capaz de reconstruir
qué información se ve, cuál se identifica y cuál debe permanecer como
memoria.

La cadena causal del atlas queda, por tanto, fijada:
\[
 \boxed{
 \begin{gathered}
 \text{extensión superviviente}
 \longrightarrow\text{ruta enriquecida}
 \longrightarrow\text{registro}
 \longrightarrow\text{firma}\\
 \longrightarrow\text{carácter}
 \longrightarrow\text{lectura dimensional}
 \end{gathered}}
 \tag{A19}\label{eq:cadena-rectora-atlas-rev6}
\]
Las ochenta y una celdas constituyen el dominio local; las cincuenta y seis
familias expresan coincidencias de ruta; las trece multisecciones son los
tipos internos; las cuarenta y una órbitas conservan la conjugación; y el
centro \((5,5)\) es la única excepción puntual.  La masa aparece después
como lectura de una persistencia ya construida.  El atlas no es una tabla
abreviada de resultados: es la estructura que explica por qué puede existir
una tabla.

\endgroup
\begingroup
\let\chapter\section
\let\section\subsection
\let\subsection\subsubsection
\let\subsubsection\paragraph
\let\clearpage\relax
\section{De las semillas APP al registro de mezcla y masa}
\label{sec:semillas-registro-mezcla-autonoma}
\index{semilla!factorización suma--producto}
\index{mezcla!procedencia en las semillas}

La ley física recibe un dominio que ya ha conservado la incompatibilidad de
las dos evaluaciones APP. Si
\(\mathcal S_+\) y \(\mathcal S_\times\) denotan las semillas orientadas de
las hojas aditiva y multiplicativa, entonces
\begin{equation}
 \Omega_{\rm seed}=\mathcal S_+\times\mathcal S_\times,
 \qquad
 |\Omega_{\rm seed}|=324^2=104,976.
 \label{eq:dominio-semillas-md-autonoma}
\end{equation}
La factorización exacta del emisor produce
\begin{equation}
 104,976
 \longrightarrow18\times26=468\ U_6
 \longrightarrow243\ w_6\subset\mathbb F_3^6,
 \qquad |\mathbb F_3^6|=729.
 \label{eq:embudo-semillas-md-autonoma}
\end{equation}
Los cuatro cardinales cuentan objetos distintos: condiciones iniciales de
dos hojas, emisiones decimales, palabras ternarias publicadas y ambiente
algebraico. Cada semilla conserva, además de su emisión, multiplicidad,
diagonal, hoja, órbita, firma de emisión, residuos módulo tres y módulo
nueve, giros, cambios de hoja, extremo, desplazamientos y conteos
cardinales. Esos datos constituyen el vocabulario del registro posterior.

La factorización organiza seis microfibras de cardinal \(1008\), seis planos
de Fano de fase, treinta y seis fibras de cardinal \(28\) y la banda
\(1944=27\cdot72\), diagonalizada por \((\mathbb Z/3)^3\). Al promediar la
palabra nonádica del TPK, los proyectores condicionales de las dos hojas
definen el núcleo canónico
\begin{equation}
 K_{\rm cat}=\frac13(I+E_++E_\times),
 \qquad
 \operatorname{Spec}(K_{\rm cat})
 =\left\{1^{[1]},
 \left(\frac23\right)^{[42]},
 \left(\frac13\right)^{[425]}\right\}.
 \label{eq:kernel-canonico-semillas-md-autonoma}
\end{equation}
La fase TPK selecciona la renovación activa y conserva las nueve bandas
espaciales; el promedio~\eqref{eq:kernel-canonico-semillas-md-autonoma}
publica la sombra estocástica del calendario completo.

\subsection{Bidegraduación de las semillas y \(6\) planos de incidencia}

En cada una de las seis microfibras de cardinalidad \(1008\), el grafo
bipartito conserva \(54\) vértices aditivos y \(56\) multiplicativos. Los
primeros se descomponen en \(36\) vértices de grado \(24\) y \(18\) de
grado \(8\); los segundos tienen grado \(18\). El recuento de incidencias
coincide por ambas hojas:
\begin{equation}
 1008=36\cdot24+18\cdot8=56\cdot18.
 \label{eq:bidegraduacion-microfibra-masas-autonoma}
\end{equation}
Esta igualdad fija la bidegraduación del emisor antes de proyectarlo sobre
las familias de ruta.

Las seis palabras
\begin{equation}
 001112,\quad010211,\quad100121,\quad112001,\quad121100,\quad211010
 \label{eq:seis-palabras-microfibra-masas-autonoma}
\end{equation}
forman una órbita de \(D_3\simeq S_3\) y poseen el perfil
\(432+4\cdot144=1008\). El cociente \(1008=36\cdot28\) organiza siete
bloques
\[
 Q=\{B_3,B_6,B_9,P_1,P_2,P_3,f\},
\]
donde \((B_3,B_6,B_9)\) son las tétradas transversales y
\((P_1,P_2,P_3,f)\), las cuatro tétradas laterales. Para una de las seis
palabras \(a\), sea
\[
 \mathcal R_a
 :=\{(p,t)\in\mathcal S_+\times\mathcal S_\times:
       U_6(p,t)\text{ publica la palabra }a\}
\]
el conjunto de sus mil ocho rutas, y pongamos
\(R_9:=\mathbb Z/9\mathbb Z\). Fijada la carta orientada, la proyección
de fase tiene tipo
\begin{equation}
 q_\beta:\mathcal R_a\longrightarrow\binom{R_9}{2},
 \qquad
 q_\beta(p,t)
 =\{x-5\beta(d_t),x+5\beta(d_t)\}\subset\mathbb Z/9\mathbb Z,
 \qquad x=5(i_p-j_p),
 \label{eq:proyeccion-fano-semillas-masas-autonoma}
\end{equation}
con \((\beta(N),\beta(E),\beta(S),\beta(O))=(1,2,3,4)\). Su codominio es
el conjunto de los treinta y seis pares no ordenados de \(R_9\). Cada
\(X\in Q\) es una tétrada de órbitas;
su soporte de rutas y su corte sobre una fibra se definen por
\[
 \operatorname{Supp}_{\rm rutas}(X)
 :=\bigcup_{\mathcal O\in X}\mathcal O,
 \qquad
 X_b:=\operatorname{Supp}_{\rm rutas}(X)\cap q_\beta^{-1}(b),
 \qquad b\in\binom{R_9}{2}.
\]
En el canal multiplicativo, el observable fuente es
\begin{equation}
 \rho:\mathcal S_\times\longrightarrow R_9,
 \qquad
 \rho(i,j,d)=(i+1)(j+1)\pmod9.
 \label{eq:residuo-producto-fano-semillas-rev3}
\end{equation}
Si \(\pi_\times:\mathcal R_a\to\mathcal S_\times\),
\(\pi_\times(p,t)=t\), es la segunda proyección, el observable que actúa
sobre rutas es
\begin{equation}
 \rho_\times:=\rho\circ\pi_\times:\mathcal R_a\longrightarrow R_9,
 \qquad
 \rho_\times(p,t)=(i_t+1)(j_t+1)\pmod9.
 \label{eq:residuo-producto-rutas-fano-semillas-rev3}
\end{equation}
En lo que sigue, \(\rho(X)\) abrevia el multiconjunto de valores de
\(\rho_\times\) sobre \(X_b\), que es independiente de la fibra \(b\). Sus
perfiles sobre las tétradas son
\[
 \rho(B_3)=3^4,
 \qquad
 \rho(B_6)=6^4,
 \qquad
 \rho(B_9)=0^4,
\]
y, para las tres tétradas laterales,
\begin{equation}
 \begin{array}{c|c|c}
 r&\text{firma activa de }P_r&R(P_r)\\
 \hline
 1&933&\{5,7\}\\
 2&393&\{4,8\}\\
 3&339&\{1,2\}.
 \end{array}
 \label{eq:perfiles-laterales-fano-semillas-rev3}
\end{equation}
Cada residuo de \(R(P_r)\) aparece dos veces entre las cuatro rutas de
\((P_r)_b\).

\begin{theorem}[Morfismo de incidencia de las semillas y seis planos de Fano]
\label{thm:semillas-seis-fano-fase-rev3}
Para \(\delta\in R_9\), \(h\in\{3,6,0\}\) ---correspondiente,
respectivamente, a \(B_3,B_6,B_9\)--- y \(r\in\{1,2,3\}\), sea
\begin{equation}
 N_\delta(B_h,P_r)
 :=\sum_{b\in\binom{R_9}{2}}
 \#\left\{(x,y)\in(B_h)_b\times(P_r)_b:
 \rho_\times(y)-\rho_\times(x)=\delta\right\}.
 \label{eq:conteo-incidencia-fano-semillas-rev3}
\end{equation}
Entonces
\begin{equation}
 N_\delta(B_h,P_r)=
 \begin{cases}
 288,&h+\delta\in R(P_r),\\
 0,&h+\delta\notin R(P_r).
 \end{cases}
 \label{eq:conteo-288-cero-fano-semillas-rev3}
\end{equation}
Si \(\delta\in R_9^\times\), para cada \(B_h\) existe un único \(P_r\)
con conteo \(288\); por tanto se obtiene una biyección
\[
 \theta_\delta:\{B_3,B_6,B_9\}
 \xrightarrow{\sim}\{P_1,P_2,P_3\}.
\]
Escribiendo la imagen de \((B_3,B_6,B_9)\) mediante sus subíndices, las seis
biyecciones son
\begin{equation}
 \begin{array}{c|c@{\qquad}c|c}
 \delta&\theta_\delta&\delta&\theta_\delta\\
 \hline
 1&(2,1,3)&5&(2,3,1)\\
 2&(1,2,3)&7&(3,2,1)\\
 4&(1,3,2)&8&(3,1,2).
 \end{array}
 \label{eq:tabla-theta-delta-fano-semillas-rev3}
\end{equation}
Sea \(P=\{P_1,P_2,P_3\}\) y, para \(p\in P\), sea
\[
 \mu(p):=\bigl\{\{f,p\},P\setminus\{p\}\bigr\}
\]
el emparejamiento perfecto asociado sobre las cuatro tétradas laterales.
Entonces
\begin{equation}
 \mathcal D_\delta
 :=\bigl\{\{B_3,B_6,B_9\}\bigr\}
 \cup
 \bigl\{\{B,q,q'\}:B\in\{B_3,B_6,B_9\},
       \{q,q'\}\in\mu(\theta_\delta(B))\bigr\}
 \label{eq:sistema-fano-semillas-rev3}
\end{equation}
es un sistema de Steiner \(\operatorname{S}(2,3,7)\). Las seis unidades
producen los seis planos de Fano sobre estos siete bloques que contienen la
línea transversal \(\{B_3,B_6,B_9\}\).
\end{theorem}

\begin{proof}
Fijemos una fibra \(b\). En \((B_h)_b\) existen cuatro elecciones con
residuo \(h\). Si \(h+\delta\in R(P_r)\), el perfil
\eqref{eq:perfiles-laterales-fano-semillas-rev3} proporciona exactamente dos
elecciones en \((P_r)_b\) con ese residuo; si no pertenece al soporte, no
proporciona ninguna. Como hay treinta y seis fibras,
\(N_\delta=36\cdot4\cdot2=288\) en el primer caso y cero en el segundo.
Los tres soportes laterales particionan \(R_9^\times\), de modo que una unidad
\(\delta\) determina un único lateral para cada bloque transversal. La lectura
directa de esos soportes da la tabla
\eqref{eq:tabla-theta-delta-fano-semillas-rev3} y agota las \(3!=6\)
biyecciones.

La línea transversal cubre las tres parejas entre \(B_3,B_6,B_9\). Para un
\(B\) fijo, las dos aristas del emparejamiento \(\mu(\theta_\delta(B))\)
cubren una vez los cuatro puntos laterales. Los tres emparejamientos son
distintos y forman la factorización de \(K_4\) en sus tres emparejamientos
perfectos; por tanto cada
pareja lateral aparece exactamente una vez. Las siete ternas contienen así
cada una de las veintiuna parejas de puntos una vez, que es la propiedad
\(\operatorname{S}(2,3,7)\). Recíprocamente, todo plano de Fano que contenga
la línea transversal asigna a cada \(B\) un emparejamiento perfecto distinto
de los otros dos y recupera una de las seis biyecciones anteriores.
\end{proof}

La correspondencia de conjuntos \(\delta\mapsto\theta_\delta\) no es una
acción ni un homomorfismo: \(R_9^\times\simeq C_6\) es abeliano bajo el
producto, mientras \(S_3\) no lo es, y las unidades no forman un subgrupo
aditivo. El desplazamiento \(\delta\) indexa seis traslaciones de los perfiles
del producto.

El mismo cálculo y la misma tabla se verifican en las seis microfibras. Debe
conservarse, no obstante, la diferencia entre soporte y rol: el soporte de
\(f\) coincide con uno de los soportes laterales y, por ello, el observable
\(\rho\) no recupera por sí solo la descomposición \(P\sqcup\{f\}\). Esa
distinción exige la firma aditiva ya declarada en el estado de semilla.

\begin{statusbox}[title={Estatuto y alcance de la incidencia de Fano}]
El \cref{thm:semillas-seis-fano-fase-rev3} es un teorema finito una vez
fijados \(q_\beta\), la orientación y el orden digital: define seis calibres
exactos de incidencia sobre bloques. Organiza fase e incidencia en el dominio
de semillas; no selecciona partícula, masa ni ruta. Su realización física se
efectúa después mediante los morfismos del registro enriquecido.
\end{statusbox}

El paso a la física conserva la cadena
\begin{equation}
 \begin{aligned}
 \text{semilla refinada}
 &\longrightarrow(\text{firmas }+,\times;U_6;w_6;
                   \text{memoria de hoja})\\
 &\longrightarrow\text{registro CT108}
 \longrightarrow(\tau_{12},D_k,\Omega,P_6,\nu_{120},\nu_{270})\\
 &\longrightarrow\text{fibra y multisección}
 \longrightarrow(B_D,B_\Omega)
 \longrightarrow\text{transporte CKM/PMNS}\\
 &\longrightarrow\text{carácter y torres}
 \longrightarrow\text{sección de masa}.
 \end{aligned}
 \label{eq:cadena-semilla-mezcla-masa-autonoma}
\end{equation}
CKM realiza un cambio de base de los operadores quark,
\begin{equation}
 B_d^{(u)}=V_{\rm CKM}B_dV_{\rm CKM}^*,
 \qquad
 R_{\rm act}^{B_d^{(u)}}
 =V_{\rm CKM}R_{\rm act}^{B_d}V_{\rm CKM}^*,
 \label{eq:ckm-transporte-semillas-autonoma}
\end{equation}
y preserva espectro y determinante. La mezcla especifica la representación
del operador en una base de sabor; el registro conserva los autovalores y su
genealogía.

\subsection{Lector neutrínico contenido en el registro}

El sector neutro dispone de un extractor entero construido exclusivamente
con datos del registro. Para cada bloque dodecafásico \(E_m\), definimos
\begin{equation}
 T(m)=\sum_{t\in E_m}(-1)^{\kappa(t)}\varepsilon(t)
       \mathbf1_{\{d(t)=9\}},
 \qquad
 s_C^{(m)}(\nu)=T(m)-T(m+6),
 \label{eq:lector-neutrino-ledger-autonoma}
\end{equation}
y
\begin{equation}
 \tau_m^\nu=\operatorname{sgn}s_C^{(m)}(\nu).
 \label{eq:trit-neutrino-ledger-autonoma}
\end{equation}
Los cuatro triángulos de salto cuatro producen
\begin{equation}
 \nu_{120}=\sum_{j=1}^{4}
 \operatorname{sgn}\!\left(\sum_{m\in\mathcal T_j}\tau_m^\nu\right),
 \qquad
 \nu_{270}=\sum_{m=1}^{12}\tau_m^\nu\tau_{m+3}^\nu.
 \label{eq:invariantes-neutrino-ledger-autonoma}
\end{equation}
La inversión bidireccional transforma
\(\nu_{120}\mapsto-\nu_{120}\) y conserva \(\nu_{270}\). Esta paridad
define la arquitectura histórica de un lector independiente de magnitudes
metrológicas empleadas como blanco. Su aplicación
literal al registro vivo \(81\times108\), tomando como
\(\varepsilon(t)\) el único bit global de signo conservado por esa
proyección, produce
\[
 s_C^{(m)}(\nu)=0,
 \qquad \tau_m^\nu=\nu_{120}=\nu_{270}=0
\]
en las ochenta y una posiciones. Los eventos antipodales han quedado
identificados por la proyección del registro. El resultado local se conserva
como lector recuperado y, a la vez, como test de pérdida de información: el
extractor neutro de
\cref{sec:extractor-neutro-z0-torsion-rev2} actúa antes de esa
identificación y lee el bit bidireccional local, la hoja, la torsión y la
corona antipodal. El cero obtenido después de proyectar no demueve esa
construcción; certifica exactamente qué información borra el lector global.
Los marcos
\(F_\ell,F_\nu\), el lector~\eqref{eq:invariantes-neutrino-ledger-autonoma}
y la forma agregada
\begin{equation}
 G_{ab}=\frac1{\sqrt{|L_a||N_b|}}
 \sum_{c\in L_a}\sum_{d\in N_b}K_{\rm registro}(c,d),
 \qquad U_{\rm int}=\operatorname{polar}(G),
 \label{eq:kernel-ledger-pmns-autonoma}
\end{equation}
fijan así los dos extremos ya construidos del puente PMNS.

La igualdad cardinal entre los cincuenta y seis vértices multiplicativos de
semillas y las cincuenta y seis familias CT108 ha sido sometida a un primer
control independiente de magnitudes metrológicas empleadas como blanco. El
mapa de Hadamard natural sobre una sola semilla alcanza
sólo veintinueve familias; al recorrer las simetrías diédricas, signos y
traslaciones naturales, la cobertura máxima es treinta y cinco. Por tanto el
entrelazador actúa sobre el par de semillas enriquecidas y conserva hoja,
acarreo y memoria. La matriz de incidencia que compare emisiones compartidas
en ese dominio doble deberá tener rango cincuenta y seis, ser natural bajo
orientación y no factorizar por contadores globales. Este criterio fija el
dominio exacto de la siguiente construcción.

\endgroup
\begingroup
\let\chapter\section
\let\section\subsection
\let\subsection\subsubsection
\let\subsubsection\paragraph
\let\clearpage\relax
\chapter{Álgebra qutrit de color y conexión finita de calibre}
\label{chap:algebra-color-qutrit}
\index{color!qutrit}\index{operadores de Weyl}
\index{grupo especial unitario}\index{acción de Wilson}

El funcional residual \(\kappa_3\) del
\cref{chap:ocupacion-estadistica} suministra tres clases afines; el atlas
posterior realiza su balance \(12+12+12\) en la
\cref{prop:color-atlas}. Una terna de etiquetas no es todavía una
representación. El paso algebraico requiere declarar la carta qutrit, su
orientación y su carácter de fase. Bajo esas hipótesis, la representación y
su álgebra infinitesimal se obtienen exactamente, sin utilizar masas ni datos
cromodinámicos.

\section{Par de Weyl y álgebra de traza \(0\)}
\label{sec:weyl-color-qutrit}

Sea
\[
 V_{\rm col}=\mathbb C^3,
 \qquad
 \omega=e^{2\pi i/3},
\]
con base ordenada \((e_0,e_1,e_2)\). Definimos
\[
 Xe_j=e_{j+1\bmod3},
 \qquad
 Ze_j=\omega^j e_j.
\]
Se cumplen
\[
 X^3=Z^3=I,
 \qquad
 ZX=\omega XZ.
\]

\begin{theorem}[Base de Weyl y álgebra compleja de color]
\label{thm:weyl-sl3-color}
Para
\[
 W_{ab}:=X^aZ^b,
 \qquad
 (a,b)\in(\mathbb Z/3\mathbb Z)^2,
\]
los nueve operadores \(W_{ab}\) forman una base ortogonal de
\(\operatorname{End}_{\mathbb C}(V_{\rm col})\) respecto del producto de
Hilbert--Schmidt. Los ocho operadores no identitarios poseen traza cero y
\[
 \operatorname{span}_{\mathbb C}
 \{W_{ab}:(a,b)\ne(0,0)\}
 =\mathfrak{sl}_3(\mathbb C).
\]
\end{theorem}

\begin{proof}
La relación de Weyl permite reducir cualquier producto a la forma
\(X^aZ^b\). Un cálculo directo da
\[
 \operatorname{tr}(W_{ab}^{\dagger}W_{cd})
 =3\,\delta_{ac}\delta_{bd}.
\]
Por tanto, los nueve operadores son linealmente independientes; como
\(\dim_{\mathbb C}\operatorname{End}_{\mathbb C}(V_{\rm col})=9\), forman una
base. La traza de \(X^aZ^b\) se anula salvo para
\((a,b)=(0,0)\). Los ocho restantes forman así una base del subespacio de
traza cero, cuya dimensión compleja es ocho.
\end{proof}

\begin{corollary}[Forma real compacta]
\label{cor:su3-color-qutrit}
La involución
\[
 \tau(A):=-A^\dagger
\]
preserva \(\mathfrak{sl}_3(\mathbb C)\), y su conjunto de puntos fijos es
\[
 \mathfrak{sl}_3(\mathbb C)^\tau
 =\{A:A^\dagger=-A,\ \operatorname{tr}A=0\}
 =\mathfrak{su}(3).
\]
Esta forma real posee dimensión ocho.
\end{corollary}

\begin{proof}
Se tiene
\[
 W_{ab}^{\dagger}=\omega^{ab}W_{-a,-b}.
\]
Los ocho índices no nulos forman cuatro pares adjuntos. Para
\[
 \mathcal R=\{(1,0),(0,1),(1,1),(1,2)\},
\]
los operadores
\[
 A_u=W_u-W_u^\dagger,
 \qquad
 B_u=i(W_u+W_u^\dagger),
 \qquad u\in\mathcal R,
\]
son antihermíticos, de traza cero y linealmente independientes sobre
\(\mathbb R\). Constituyen, por tanto, una base real de
\(\mathfrak{su}(3)\).
\end{proof}

La aplicación residual
\[
 \kappa_{\rm col}(r,c)
 =\mathbb C e_{\,r-c\bmod3}
\]
asocia las tres clases del atlas con las tres rectas propias de \(Z\). El
operador \(X\) las permuta cíclicamente y \(Z\) registra su fase. Esta
asignación es exacta una vez fijados
\[
 V_{\rm col}=\mathbb C[\mathbb F_3],
 \quad
 \text{la base delta},
 \quad
 j\longmapsto j+1,
 \quad
 j\longmapsto\omega^j.
\]
Los cambios afines admisibles de la carta producen representaciones
unitariamente conjugadas. La construcción es, por ello, natural hasta
conjugación unitaria; no selecciona por sí sola los nombres físicos de los
colores ni una realización QCD.

\section{Conexión y acción sobre un complejo celular finito}
\label{sec:gauge-finito-color}

\begin{construction}[Teoría gauge finita relativa al complejo]
\label{cons:gauge-finito-color}
Sea \(\mathcal K=(V,E,F)\) una celularización finita orientada compatible con
el grafo de rutas. A cada arista orientada \(e\) se asigna
\[
 U_e\in SU(3),
 \qquad
 U_{\bar e}=U_e^{-1}.
\]
Una transformación \(g:V\to SU(3)\) actúa por
\[
 U_e\longmapsto U_e^g
 =g_{t(e)}U_eg_{s(e)}^{-1}.
\]
Para cada cara \(f\), el producto ordenado
\[
 U_f=\prod_{e\subset\partial f}U_e^{\epsilon(f,e)}
\]
define su holonomía. Elegido un vértice base \(v_f\), se tiene
\[
 U_f^g=g_{v_f}U_fg_{v_f}^{-1}.
\]
\end{construction}

\begin{proposition}[Invariancia gauge y transporte covariante]
\label{prop:wilson-color-finito}
Para \(\beta\ge0\), la acción
\[
 S_{\rm col}^{(\mathcal K,\beta)}[U]
 =\frac{\beta}{3}\sum_{f\in F}
 \left(3-\operatorname{Re}\operatorname{tr}U_f\right)
\]
es gauge-invariante y no negativa. Si
\(\psi_v\in V_{\rm col}\) transforma como
\(\psi_v\mapsto g_v\psi_v\), entonces
\[
 \nabla_e\psi:=U_e\psi_{s(e)}-\psi_{t(e)}
\]
satisface
\[
 \nabla_e^g\psi^g=g_{t(e)}\nabla_e\psi.
\]
\end{proposition}

\begin{proof}
La holonomía cambia por conjugación y la traza es cíclica; de ahí la
invariancia. Como cada \(U_f\) es unitario de dimensión tres,
\(\operatorname{Re}\operatorname{tr}U_f\le3\), lo que prueba la no
negatividad. La covariancia de \(\nabla_e\) resulta sustituyendo las dos leyes
de transformación.
\end{proof}

\section{Positividad de núcleos de clase y descomposición armónica}

La conexión finita precedente no implica por sí sola positividad por
reflexión ni un límite de Yang--Mills. Antes de formular tales transferencias
conviene aislar el control clásico que sí está cerrado. Sea \(G\) un grupo
compacto y \(\widehat G\) su dual unitario. Para cada
\(\rho\in\widehat G\), sean \(d_\rho\) su dimensión y
\(c_\rho(a)\ge0\) coeficientes con decaimiento suficiente, respecto del
crecimiento de \(d_\rho\), para que la serie siguiente converja
absolutamente.

\begin{definition}[Núcleo de clase]
\label{pf84:YM-07}
Se define
\[
 K_a(g,h):=
 \sum_{\rho\in\widehat G}
 c_\rho(a)\,
 \Tr\!\bigl(\rho(g)\rho(h)^*\bigr),
 \qquad g,h\in G.
\]
La elección
\[
 c_\rho(t)=d_\rho e^{-tC_2(\rho)},
 \qquad t>0,
\]
proporciona el ejemplo del núcleo de calor. La no negatividad de los
coeficientes no reemplaza la hipótesis de convergencia absoluta.
\end{definition}

\begin{theorem}[Positividad de Peter--Weyl]
\label{pf84:YM-08}
Para \(g_1,\dots,g_n\in G\) y \(z_1,\dots,z_n\in\CC\),
\[
 \sum_{i,j=1}^n
 \overline z_i z_jK_a(g_i,g_j)\ge0.
\]
\end{theorem}

\begin{proof}
La unitariedad de las representaciones y la convergencia absoluta permiten
reordenar la suma:
\[
 \sum_{i,j}\overline z_i z_jK_a(g_i,g_j)
 =
 \sum_{\rho\in\widehat G}
 c_\rho(a)
 \left\lVert
 \sum_j\overline z_j\rho(g_j)
 \right\rVert_{\mathrm{HS}}^2\ge0.
\]
Un coeficiente espectral negativo puede destruir esta positividad de Gram;
por ello el signo forma parte de las hipótesis.
\end{proof}

\begin{lemma}[Control abeliano de Poisson]
\label{pf84:YM-09}
Para \(G=U(1)\), \(\theta\in\RR\) y \(|a|<1\),
\[
 \sum_{n\in\ZZ}a^{|n|}e^{in\theta}
 =\frac{1-a^2}{1-2a\cos\theta+a^2}.
\]
Si se exigen además coeficientes de Fourier no negativos, se restringe
\(a\) a \(0\le a<1\).
\end{lemma}

\begin{proof}
La igualdad resulta de sumar por separado las dos series geométricas de
índices positivos y negativos y añadir el modo \(n=0\). Para \(a<0\), la
identidad escalar sigue siendo válida, pero ya no constituye un ejemplo de
coeficientes Peter--Weyl no negativos.
\end{proof}

\begin{remark}
Los \cref{pf84:YM-07,pf84:YM-08,pf84:YM-09} son definiciones y resultados
clásicos de análisis armónico. Proporcionan un control positivo para
núcleos de clase; no demuestran positividad Osterwalder--Schrader, brecha de
masa, confinamiento ni límite continuo, y no promueven la teoría gauge
finita de este capítulo a cromodinámica cuántica.
\end{remark}

\begin{exactbox}[title={Cierre algebraico del color}]
Los \cref{thm:weyl-sl3-color,cor:su3-color-qutrit} construyen exactamente la
cadena
\[
\text{residuo ternario}\longrightarrow
\text{qutrit}\longrightarrow
\mathfrak{sl}_3(\mathbb C)\longrightarrow\mathfrak{su}(3).
\]
La \cref{cons:gauge-finito-color,prop:wilson-color-finito} prolonga esa
cadena a una conexión y una acción gauge sobre la celularización finita
\(\mathcal K\).  Éste es el objeto demostrado en el capítulo.  La
comparación con la cromodinámica física se efectúa después y no redefine el
qutrit, el balance \(12+12+12\) ni la representación construida.
\end{exactbox}

\endgroup
\begingroup
\let\chapter\section
\let\section\subsection
\let\subsection\subsubsection
\let\subsubsection\paragraph
\let\clearpage\relax
% Vista editorial de corpus/source/masas/02_ley_general.tex, líneas 363--491: cuerpo exacto salvo el retítulo académico del epígrafe sobre el ciclo de 108 estados. Las líneas 1--263 ya residen exactamente en 72_rev2_nucleo_ley_i.tex y las líneas 264--362 en 72_rev2_algoritmo_universal.tex.
\chapter{Ampliación sistemática de la teoría de masas}
\Needspace{7\baselineskip}
\section{Fundamento unificado de la masa}
\Needspace{7\baselineskip}
\subsection{El objeto anterior a la magnitud}
La teoría comienza antes de la cifra. Una masa físicamente individualizada no es un número real acompañado por el nombre de una partícula, sino el término final de una biografía matemática. APP aporta el espacio finito de incidencias y las dos operaciones correlativas; el TRIT conserva la orientación temporal de cada transición; el TPK transporta fase, hoja, acarreo, frontera, ruta, memoria y supervivencia. La Mecánica Dimensional convierte después esa biografía en una sección de la recta de masa. El orden completo es

\[
\mathrm{APP}\longrightarrow\mathrm{TRIT}\longrightarrow\mathrm{TPK}
\longrightarrow\Gamma^{\rm surv}
\longrightarrow\gamma_{\rm obs}
\longrightarrow\mathscr L_\gamma
\longrightarrow\sigma(\gamma)
\longrightarrow\chi_{\rm mass}
\longrightarrow\mathcal T_{90,120}
\longrightarrow\mathcal L_M.
\]
Cada flecha retiene información que la siguiente utiliza. La celda no reemplaza la ruta; la ruta no reemplaza el registro; una firma no reemplaza su registro genealógico; un carácter no reemplaza la sección dimensional sobre la que actúa. Esta separación explica por qué dos partículas pueden compartir carga o espín y, sin embargo, poseer masas diferentes: su biografía enriquecida no es la misma. También explica por qué partícula y antipartícula comparten masa: la involución cambia la orientación impar y conserva la memoria par que desciende al cociente de masa.

APP consta de nueve por nueve posiciones publicadas. Cada posición admite cuatro orientaciones iniciales, de modo que el dominio orientado contiene

\[
81\cdot4=324
\]
rutas candidatas. La clasificación interna produce 56 familias de ruta y 13 multisecciones. Estas cardinalidades son simultáneas: 81 cuenta celdas, 324 cuenta orientaciones, 56 cuenta cocientes por historia y 13 cuenta tipos operatorios gruesos. Ninguna tabla de partículas sustituye a esos dominios.

\Needspace{7\baselineskip}
\subsection{Unidad discreta del continuo que sostiene la masa}
La recta de masa no se adosa a cinco teorías independientes del continuo. Procede del único estado APP--TRIT--TPK del que emergen correlativamente los aspectos espectral, solenoidal, de calibre, coherente y determinantal, designados en la notación de construcción por \(sp,sol,gau,coh,det\). Esos símbolos son discriminantes de lectura, no subdominios previos. Si \(\mathcal C_{\rm disc}\) denota la estructura común, la flecha generadora tiene la forma

\[
\mathfrak E_{\rm cont}:\mathcal C_{\rm disc}
\longrightarrow
\bigl(\mathcal C_\infty;
\mathcal O_{sp},\mathcal O_{sol},\mathcal O_{gau},
\mathcal O_{coh},\mathcal O_{det}\bigr),
\]
donde las cinco observaciones pertenecen a un mismo objeto límite, comparten fibras, orientación, acarreo, memoria, frontera y lectores, y se transportan por los mismos cuadrados de refinamiento. La masa utiliza su lectura determinantal para la escala de acción, su lectura espectral para el carácter y las torres, su lectura de calibre para los sectores nulos y cargados, y sus lecturas coherente y solenoidal para la propagación y la composición. Ninguna de ellas se introduce suponiendo de antemano que el estado pertenece a una clase física clásica.

Esta unidad es una condición de tipado: separar uno de esos aspectos y conservar los otros produciría un objeto distinto, incapaz de sostener a la vez ruta, acción, carga, amplitud y masa. La presente teoría no vuelve a demostrar el continuo completo; incorpora su construcción unitaria como fundamento y verifica que ninguna de las inserciones de masas lo disgrega.

\Needspace{7\baselineskip}
\subsection{Escala absoluta y cocientes}
El lector determinantal local tiene forma normal de Smith

\[
\operatorname{SNF}(H_4)=\operatorname{diag}(1,2,2,4),
\]
de donde

\[
|\operatorname{coker}H_4|=16,
\qquad
\Sigma_H=\varphi_{\rm HMT}\alpha_{\rm HMT}^{16},
\qquad
\operatorname{ord}_{10}(\Sigma_H)=-34.
\]
La década se obtiene antes de la carta SI. Las secciones pre y retorno de la recta de acción son

\[
\hbar_{\rm pre}=H_5\,10^{-34}u_S,
\qquad
\hbar_{\rm ret}=\eta_{\rm ret}^{(\deg)}\,10^{-34}u_S,
\]
y su cociente bidireccional es

\[
R_{\rm act}=\frac{H_5}{\eta_{\rm ret}^{(\deg)}}.
\]
En la rama de reloj elegimos explícitamente como sección de acción la sección de retorno, \(\hbar_{\rm int}:=\hbar_{\rm ret}\). Un reloj declarado \(t_0\) proporciona

\[
\omega_0=\frac{2\pi}{108t_0},
\qquad
\varepsilon_{\rm clk}=\hbar_{\rm ret}\omega_0,
\qquad
m_{\rm clk}=\varepsilon_{\rm clk}c_{\rm int}^{-2}.
\]
Adoptar \(u_M^{\rm clk}:=m_{\rm clk}\) como origen de la carta es una normalización explícita. No identifica automáticamente esta sección con una base energética o másica previamente declarada. El carácter produce primero el operador basal

\[
\widehat M_X^{(0)}
=\mathfrak e_0m_{\rm clk}\,
\chi_{\rm full}(\sigma_X-\sigma_e,\eta_X-\eta_e),
\]
y el lector bidireccional actúa después. Para cada presentación \(r\in\{D,\Omega\}\), sea \(\widehat B_X^r=(\widehat B_X^r)^*\) sobre la fibra finita; en una realización infinita se exige además un dominio común denso para el cálculo funcional. Entonces

\[
\widehat M_X^{\beta,r}
=R_{\rm act}^{\widehat B_X^r/2}
\widehat M_X^{(0)}
R_{\rm act}^{\widehat B_X^r/2}.
\]
En una fibra escalar, \(\widehat B_X^r=\kappa_XI\), de modo que

\[
m_X^\beta=m_X^{(0)}R_{\rm act}^{\kappa_X}.
\]
Si las coordenadas de torre ya se expresan relativamente al electrón, se omite \(-\eta_e\). Al formar un cociente, la sección dimensional común se cancela, pero no la diferencia de lectura bidireccional:

\[
\frac{m_Y^\beta}{m_X^\beta}
=\chi_{\rm full}(\sigma_Y-\sigma_X,\eta_Y-\eta_X)
R_{\rm act}^{\kappa_Y-\kappa_X}.
\]
Por tanto, \(10^{-34}\) fija la escala absoluta de acción, energía y masa; no es un corrector relativo oculto. La información fina de una especie sólo puede entrar mediante su carácter, su reloj relativo o el cociente orientado de las secciones pre y retorno.

\Needspace{7\baselineskip}
\subsection{Cuanto interno de la recta de masa inducido por el ciclo de 108 estados}
La autodualidad Planck--CT108 proporciona, en la normalización interna

\[
\hbar_{\rm int}=c_{\rm int}=t_0=1,
\]
la sección

\[
m_{108}^{\rm int}
=0.058177641733144319230789692282953757\ldots .
\]
Este número es un cuanto normalizado de la recta interna de masa. No es la masa de una partícula, no lleva por sí solo una carta MeV o GeV y no puede insertarse como fila adicional de especies. Su función es comprobar la compatibilidad entre la carta de acción, el reloj CT108 y la autodualidad de la recta de masa. La transición a una masa nominal exige todavía la ruta, la firma, el operador y la sección dimensional de la especie.

\Needspace{7\baselineskip}
\subsection{Existencia de la salida HMT–MD para extensiones supervivientes y criterio de reducción escalar por rango de fibra}
Fijada una extensión superviviente, un lector celular, una firma entera, una familia de coeficientes de torre absolutamente convergente y una sección dimensional de masa, existe una salida HMT--MD bien tipada. Si la fibra posee un único punto fijo compatible con la involución, la salida se reduce a un escalar canónico. Si la fibra tiene rango mayor, la salida primaria es el operador espectral sobre la fibra; una masa escalar física requiere el funcional de acoplamiento que individualiza la representación observada.

Este teorema no degrada las evaluaciones escalares recuperadas. Distingue tres afirmaciones: el operador interno está construido; la evaluación de una firma declarada es exacta; la determinación prospectiva de esa firma es un teorema adicional cuando no se deduce del propio lector. La comparación experimental se efectúa únicamente después de fijar cuál de esas tres afirmaciones se publica.

\Needspace{7\baselineskip}

\endgroup

\section{Algoritmo universal de realización material e igualador espectral de los lectores \(K_D\) y \(K_\Omega\)}
La construcción termina con el algoritmo que recibe únicamente datos
genealógicos tipados, forma la fibra de realización y exige el equalizador de
los lectores antes de publicar una salida dimensional.

\begingroup
\let\chapter\section
\let\section\subsection
\let\subsection\subsubsection
\let\subsubsection\paragraph
\let\clearpage\relax
\section{Algoritmo universal de la Ley General de Masas}
\Needspace{7\baselineskip}
\subsection{Dominio del lector celular: extensión superviviente, ruta, registro y firma}
La entrada contiene únicamente:

\[
 (\mathrm{APP},\mathrm{TRIT},\mathrm{TPK},
 \mathcal R_{324},\mathcal L_{108},
 \text{representación},\text{acoplamiento},
 \text{carta dimensional común}).
\]
No contiene masa observada, residual, coeficiente elegido por aproximación ni tabla externa.

\Needspace{7\baselineskip}
\subsection{Construcción functorial de la realización material: fibra de rutas, entrelazador, proyector de coincidencia, operador espectral y escalarización}
Para cada estado físico \(X\):

\begin{enumerate}
\item construir su descriptor de carga, generación, sabor, color, espín, estadística,
orientación y composición;

\item formar la fibra de rutas compatibles;
\item calcular el espacio de entrelazadores
\end{enumerate}
   \[
   \mathscr I_X=
   \operatorname{Hom}_{G_X}(V_X,\ell^2(F_X));
   \]
\begin{enumerate}
\item si \(\dim\mathscr I_X=0\), descartar esa identificación;
\item construir el proyector mediante promedio de Reynolds y factor polar;
\item reconstruir los registros genealógicos de todas las rutas del soporte;
\item calcular firma, \(K_D\), \(K_\Omega\) y
\(\mathcal F_{\rm tower}^{2w}\);

\item aplicar el puntero de lectura del acoplamiento, el equalizador o la medida
espectral conjunta;

\item construir \(\widehat M_X\);
\item reducir a escalar sólo si la varianza es cero o existe un polo aislado;
\item para compuestos, aplicar antes \(\mathfrak C_{12}\);
\item para anchuras, calcular el semigrupo de supervivencia;
\item multiplicar por la sección dimensional común;
\item congelar la salida;
\item abrir la comparación externa.
\end{enumerate}
\Needspace{7\baselineskip}
\subsection{Existencia de la realización por fibras}
La multiplicidad de una representación \(\pi_X\) dentro de una fibra se calcula por

\[
 \dim\mathscr I_X=
 \frac1{|G_X|}
 \sum_{g\in G_X}
 \overline{\chi_{\pi_X}(g)}
 |\operatorname{Fix}_{F_X}(g)|.
\]
Si es positiva, existe un entrelazador. Si es mayor que uno, el resultado es un haz de realizaciones y el acoplamiento debe seleccionar un estado o conservar la medida mixta. Esta fórmula reemplaza una búsqueda nominal por una condición algebraica finita sobre las 324 rutas.

\Needspace{7\baselineskip}
\subsection{Igualador espectral de los lectores \(K_D\) y \(K_\Omega\) mediante \(Q_{\mathrm{eq}}=\mathbf1_{\{0\}}(\widehat L_D-\widehat L_\Omega)\)}
Sean

\[
 L_r(\Gamma)=
 \kappa_r(\Gamma)\log R_{\rm act}
 +\sum_{h\in\mathfrak S}\eta_h^r(\Gamma)s_h,
 \qquad r\in\{D,\Omega\}.
\]
El proyector de coincidencia es

\[
 Q_{\rm eq}=
 \mathbf1_{\{0\}}(\widehat L_D-\widehat L_\Omega).
\]
Una realización admite lectura común exactamente cuando

\[
 P_X\le Q_{\rm eq}.
\]
Si existe una involución del acoplamiento que intercambia ambos lectores, se usa \(L_{\rm sym}\). En otro caso, se publica su medida espectral conjunta.

\Needspace{7\baselineskip}
\subsection{Codominio del lector celular: carácter adimensional, torre y realización de masa}
La salida puede ser:

\[
 \begin{cases}
 m_X\in\mathcal L_M,&\text{masa escalar},\\
 \operatorname{Spec}(\widehat M_X),&\text{multiplete},\\
 z_X=M_X-i\Gamma_X/2,&\text{resonancia},\\
 0\in\mathcal L_M,&\text{sección nula},\\
 (\mathcal C,F,R,\Delta E),&\text{sector topológico},\\
 (L,\operatorname{Obs}_L),&\text{sección virtual}.
 \end{cases}
\]
Esta diversidad de codominios es parte de la ley; no una excepción.

\Needspace{7\baselineskip}

\endgroup

\section{Bifurcación nula y positividad}
El sector fotónico y gluónico se separa antes de aplicar el carácter positivo.
La identidad \(R_{\rm act}^{0}=1\) no demuestra masa cero; la nulidad procede
de la carta de calibre. En las familias materiales, la salida primaria es un
operador positivo de fibra. Sólo una compresión física tipada produce un
escalar.

\section{Portal del dominio exhaustivo}
La definición y los teoremas del dominio anteceden al despliegue de sus
\hyperref[chap:atlas-genealogico-exhaustivo-324-rutas]{324 fichas completas}.
Las fichas se sitúan al final del aparato científico para conservar la
continuidad de la derivación sin ocultar ninguna ruta.
